\documentclass[11pt]{amsart}
\usepackage[margin=1.1in]{geometry}
\usepackage{amsmath,amssymb,amsthm}
\usepackage{mathtools}
\usepackage{xcolor}
\usepackage[colorlinks=true,linkcolor=blue!50!black,citecolor=blue!50!black,urlcolor=blue!50!black]{hyperref}
\usepackage[capitalise,nameinlink]{cleveref}
\usepackage{booktabs}

\newtheorem{theorem}{Theorem}[section]
\newtheorem{lemma}[theorem]{Lemma}
\newtheorem{proposition}[theorem]{Proposition}
\newtheorem{corollary}[theorem]{Corollary}
\theoremstyle{definition}
\newtheorem{definition}[theorem]{Definition}
\theoremstyle{remark}
\newtheorem{remark}[theorem]{Remark}

\crefname{lemma}{Lemma}{Lemmas}
\crefname{theorem}{Theorem}{Theorems}
\crefname{corollary}{Corollary}{Corollaries}
\crefname{proposition}{Proposition}{Propositions}
\crefname{remark}{Remark}{Remarks}

\newcommand{\E}{\mathbb{E}}
\renewcommand{\Pr}{\mathbb{P}}
\newcommand{\Inf}{\mathrm{Inf}}
\newcommand{\Var}{\mathrm{Var}}
\newcommand{\Bin}{\mathrm{Bin}}
\newcommand{\hatg}{\widehat{g}}
\newcommand{\bits}{\{0,1\}}

\title[Entropy--influence for read-$k$ DNFs and regular functions]{The Fourier entropy--influence inequality for read-$k$ DNF formulas and regular Boolean functions}
\author{Yitzchak Shmalo}
\address{Einstein Institute of Mathematics, The Hebrew University of Jerusalem, Jerusalem, Israel}
\email{yitzchak.shmalo@gmail.com}
\date{23 September 2026}

\begin{document}

\begin{abstract}
The Fourier entropy--influence conjecture of Friedgut and Kalai asserts that the spectral entropy of a Boolean
function is at most a universal constant times its total influence. We prove it with constant $15.1$ for regular
Boolean functions, a class that contains all linear threshold functions, and with constant $O(\log k)$ for read-$k$
DNF formulas and, more generally, for unions of regular functions in which every variable is relevant to at most $k$
of them; for Boolean functions of $m$ regular functions the constant is $O(\log m)$. These classes satisfy the
stronger inequality of O'Donnell and Tan, so by their composition theorem the conjecture also holds for read-once
formulas built from such functions and gates of bounded arity, including iterated majorities of any odd fan-in.
Threshold gates of any fan-in satisfy the biased form of that inequality under every product measure, so the
conjecture also holds, with an absolute constant, for read-once formulas of any depth over threshold gates of
unbounded fan-in, including iterated threshold functions with arbitrary thresholds and fan-ins.
\end{abstract}

\maketitle

\section{Introduction}

For $f:\bits^n\to\bits$ put $g=1-2f\in\{-1,1\}$ and expand $g=\sum_{A\subseteq[n]}\hatg(A)\chi_A$ with
$\chi_A(x)=(-1)^{\sum_{i\in A}x_i}$. By Parseval the numbers $P(A)=\hatg(A)^2$ form a probability distribution on
subsets of $[n]$, the \emph{spectral distribution} of $f$. Its Shannon entropy and its mean size are
\[
H(f)=\sum_{A}P(A)\log_2\frac{1}{P(A)},\qquad I(f)=\sum_A|A|\,P(A)=\sum_{i=1}^n\Inf_i(f),
\]
where $\Inf_i(f)=\Pr[f(x)\ne f(x\oplus e_i)]$ for uniform $x$. Friedgut and Kalai~\cite{FK96} conjectured that
$H(f)\le C\,I(f)$ for a universal constant $C$ (the FEI conjecture). The conjecture is open. It is known for
symmetric functions and read-once decision trees~\cite{OWZ11}, read-once formulas over gates of bounded
arity~\cite{OT13} and over $\mathrm{AND}$, $\mathrm{OR}$, $\mathrm{XOR}$ and $\mathrm{NOT}$ gates~\cite{CKLS16}, read-$k$ decision trees~\cite{WWW14,Sha18}, random DNF formulas and random linear
threshold functions~\cite{KLW10,CKKLS18}, plateaued functions~\cite{Pot26}, and $\delta$-tribes functions and functions with
separation properties~\cite{GMP26}. For every Boolean function, Kelman, Kindler, Lifshitz, Minzer and
Safra~\cite{KKLMS20} proved $H(f)\le K\sum_A|A|\log_2(|A|+1)P(A)+K\,I(f)$ with an absolute constant $K$, which
misses the conjecture by a factor logarithmic in the level. The subadditivity of entropy gives
$H(f)\le\sum_ih(\Inf_i(f))=O\bigl(I(f)+\sum_i\Inf_i(f)\log(1/\Inf_i(f))\bigr)$, with $h$ the binary
entropy~\cite{KMS12}; Han~\cite{Han25} proved the second bound by an iteration over the levels. The best known lower
bound on $C$ is
$6.5218$~\cite{Num26}, following $6.278$~\cite{OT13} and $6.4548$~\cite{Hod17}.

The conjecture implies Mansour's conjecture on the Fourier spectrum of DNF formulas~\cite{Man95,OWZ11}. Mansour's
conjecture is known for read-$k$ DNF formulas~\cite{KLW10}. The weaker Fourier min-entropy--influence inequality,
$\log_2(1/\max_A P(A))\le C\,I(f)$, was proved for regular read-$k$ DNF formulas by Shalev~\cite{Sha18} and for all
read-$k$ DNF formulas, with $C=O(\log k)$, by Arunachalam, Chakraborty, Kouck\'y, Saurabh and de Wolf~\cite{ACKSW21}.
The entropy inequality itself was open for this class; Shalev~\cite{Sha18} proposed it as a direction.

A DNF formula is a disjunction $\varphi=T_1\vee\cdots\vee T_m$ of terms, each term $T_t$ being the conjunction of
literals on a nonempty set $V_t$ of variables. It is \emph{read-$k$} if every variable lies in at most $k$ of the
sets $V_t$. No minimality is assumed.

\begin{theorem}\label{thm:readk}
Let $f$ be computed by a read-$k$ DNF formula. Then
\[
H(f)\le \min\bigl(2.89+9.9\,k,\ 30+5.6\log k\bigr)\,I(f).
\]
The same bound holds for functions computed by read-$k$ CNF formulas.
\end{theorem}

The dependence on $k$ is logarithmic, as in the min-entropy bound of~\cite{ACKSW21}. A DNF formula with $m$ terms is
read-$m$, so every such formula satisfies $H(f)\le(30+5.6\log m)\,I(f)$.

For linear threshold functions the conjecture was open: Chakraborty, Kulkarni, Lokam and Saurabh~\cite{CKLS16} proved
$H(f)=O(\sqrt n)$ for them, and~\cite{CKKLS18} treated random ones. A monotone Boolean function $f$ is
\emph{regular} (or $2$-monotone~\cite{Mur71}) if its coordinates are totally
preordered by desirability: for all $i\ne j$ either $f(x^{i\leftarrow1,j\leftarrow0})\ge f(x^{i\leftarrow0,j\leftarrow1})$
for every $x$, or the reverse inequality holds for every $x$. A function is regular if it becomes a monotone
regular function after negating some of its inputs. Every linear threshold function is regular.

\begin{theorem}\label{thm:regular}
Every regular Boolean function satisfies $H(f)\le 15.1\,I(f)$.
\end{theorem}

The proof extends to Boolean combinations of regular functions whose desirability orders are unrelated
(\cref{thm:multi}): if $h_1,\dots,h_m$ are regular and $F:\bits^m\to\bits$ is arbitrary, then
$f=F(h_1,\dots,h_m)$ satisfies $H(f)\le(16.5+8.4\ln m)\,I(f)$. In particular the conjecture holds, with constant
$O(\log m)$, for intersections of $m$ halfspaces. For unions and intersections the dependence on $m$ can be replaced
by the overlap: if $h_1,\dots,h_m$ are regular and every variable is relevant to at most $k$ of them, then
$h_1\vee\cdots\vee h_m$ and $h_1\wedge\cdots\wedge h_m$ satisfy $H(f)\le\min(3.9+21.2k,\,54+12.8\log k)\,I(f)$
(\cref{thm:union,thm:union-log}). This
contains \cref{thm:readk,thm:regular} up to the constants. The two statements combine: if every $h_l$ is a
Boolean function of at most $M$ regular functions, the constant is $54+12.8\log k+6.8\log M$ (\cref{thm:union-multi});
for example, for unions of polytopes with $M$ facets in which every variable enters at most $k$ of them, and, since
every function of $w$ variables is a function of $w$ dictators, for read-$k$ unions of arbitrary functions of at most
$w$ variables each.

O'Donnell and Tan~\cite{OT13} introduced a stronger inequality that is preserved under disjoint composition. For
the uniform measure it reads
\begin{equation}\label{eq:feiplus}
H^+(f):=\sum_{A\ne\emptyset}P(A)\log_2\frac1{P(A)}\le C\,\bigl(I(f)-\Var(f)\bigr),
\end{equation}
where $\Var(f)=1-P(\emptyset)$, and it implies $H(f)\le\max(C,1/\ln2)\,I(f)$~\cite[Proposition~2.1]{OT13}.

\begin{theorem}\label{thm:plus-intro}
Inequality \eqref{eq:feiplus} holds with $C=\min(22+28k,\,88+16.4\log k)$ for functions computed by read-$k$ DNF or CNF formulas, with
$C=60$ for regular functions, with $C=\min(24+66k,\,186+46.2\log k)$ for unions and intersections of regular functions in which every
variable is relevant to at most $k$ of them, and with $C=66+31\ln m$ for Boolean functions of $m$ regular functions.
\end{theorem}

Together with the composition theorem of~\cite{OT13}, this gives the conjecture for a class that contains read-once
formulas, read-$k$ DNF formulas and intersections of halfspaces.

\begin{corollary}\label{cor:compose-intro}
Let $F$ be a read-once formula and $f=F(f_1,\dots,f_s)$, where $f_1,\dots,f_s$ depend on pairwise disjoint sets of
variables and each $f_l$ belongs to one of the classes of \cref{thm:plus-intro}. Suppose that every gate of $F$ is an
$\mathrm{AND}$, $\mathrm{OR}$ or $\mathrm{XOR}$ gate of any fan-in, or has arity at most $a$, or is a threshold gate
$T_{t,r}(y)=\mathbf 1[y_1+\cdots+y_r\ge t]$ of any fan-in $r$, or belongs to one of these classes and has only
balanced inputs, that is, inputs equal to $1$ with probability $\frac12$. Then $H(f)\le C\,I(f)$, where $C$ is the
largest of $C_a=2^{O(a)}$, of $760$ if threshold gates occur ($66.4$ if the inputs of each threshold gate are all
equal to $1$ with the same probability), and of the constants of \cref{thm:plus-intro} for the classes that occur.
\end{corollary}

For instance, the corollary applies to unions and intersections of any number of halfspaces on pairwise disjoint
sets of variables, to majorities of intersections of halfspaces on disjoint sets of variables, and to iterated
majorities $\mathrm{MAJ}_{r_1}\circ\mathrm{MAJ}_{r_2}\circ\cdots\circ\mathrm{MAJ}_{r_d}$ with odd fan-ins of any
size, since a majority of an odd number of independent balanced bits is balanced; in these examples the constant is
$66$, independent of the number of halfspaces and of the fan-ins. Threshold gates satisfy the biased form of
\eqref{eq:feiplus} under every product measure, with constant $66.4$ when their inputs have a common bias
(\cref{thm:threshold}) and with constant $760$ in general (\cref{thm:thr-strong}). So the corollary also applies to
iterated threshold functions $T_{t_1,r_1}\circ T_{t_2,r_2}\circ\cdots\circ T_{t_d,r_d}$ with arbitrary thresholds and
fan-ins, where the bias changes from level to level, and to thresholds of conjunctions of equal width, with constant
$66.4$, and to every read-once formula over threshold gates of any fan-in, with constant $760$
(\cref{cor:thr-formula}). The constant $760$ comes from splitting at pivotal probability $\frac9{10}$: above it the
other inputs of the gate are nearly deterministic and a direct estimate applies, and below it the weak form loses a
factor $10$. The weaker inequality of \cref{prop:thr-general}, with $I_\mu$ in place
of $I_\mu-\Var_\mu$, composes at the cost of adding its constant at each level (\cref{lem:weak-compose}) and gives
$H(f)\le(38d+\log e)\,I(f)$ for formulas of depth $d$ (\cref{cor:depth}), which is smaller when $d\le19$. Gates $\mathrm{AND}$, $\mathrm{OR}$ and $\mathrm{XOR}$
of any fan-in are read-once formulas over gates of arity $2$, so~\cite{OT13} already covers them with an unspecified
constant; \cref{lem:andor} gives them the biased form of \eqref{eq:feiplus} with the sharp constant $2+4\log e$, which
makes the constant of the corollary explicit for formulas over these gates.

The proofs of the two theorems use different mechanisms. For
\cref{thm:readk}, the discrete derivative $D_ig$ is supported inside the union of at most $k$ subcubes, one for each
term containing $i$, and on each subcube it factorises exactly into the indicator of the subcube times a function
of the remaining coordinates. Consequently the spectral sample, conditioned on containing $i$, contains a uniformly
random subset of the variables of some term through $i$ (\cref{lem:witness}). This is combined with a Gibbs
inequality against a reference measure that activates terms independently, and with a path argument that bounds
the cost of the terms that are not activated (\cref{lem:path}). For the logarithmic bound only one term of each
connected piece of the cover is activated; the other terms of the piece are reached through shared variables, at a
cost of order $\log k$ each (\cref{thm:readk-log}). For \cref{thm:regular}, the pivotal set of a
coordinate, restricted to the coordinates that are at least as desirable, is an antichain; Sperner's theorem then
shows that the spectral sample, conditioned on containing $i$, contains a positive proportion of the more desirable
coordinates in a logarithmic average (\cref{lem:ranktail}), and a rank-weighted reference measure converts this into
an entropy bound. \Cref{thm:union} combines the two: the reference measure activates the functions $h_l$, and the
uniform filling of a witnessing term is replaced by a level inequality for functions supported on the pivotal set
of a rarely true $h_l$ (\cref{lem:level}); \cref{thm:union-log} adds the rooted activation of \cref{thm:readk-log}
and rank laws normalised at the size of the part of $S$ assigned to each $h_l$. For \cref{thm:plus-intro}, two lemmas valid for every Boolean function reduce
\eqref{eq:feiplus} to the sets with at least two elements (\cref{lem:single,lem:mass}); on those sets both proofs
go through with small changes (\cref{sec:plus}).

\subsection*{Notation and standard facts}
Throughout, $x$ is uniform on $\bits^n$, $\log=\log_2$, and $\ln$ is the natural logarithm. $S$ denotes a random set
with law $P$ (the spectral sample of $f$), so $H(f)=H(S)$, $I(f)=\E|S|$ and $\Pr[i\in S]=\Inf_i(f)$. We write
$\Var(f)=1-P(\emptyset)=4p(1-p)$ with $p=\Pr[f=1]$, and $h$ for the binary entropy function in bits.

For $i\in[n]$ the discrete derivative
\[
D_ig(x)=\tfrac12\bigl(g(x^{i\leftarrow0})-g(x^{i\leftarrow1})\bigr)=f(x^{i\leftarrow1})-f(x^{i\leftarrow0})
\]
does not depend on $x_i$, takes values in $\{-1,0,1\}$, and is nonzero exactly on the set $\mathcal P_i$ of points
at which $i$ is pivotal. Since $g=\E_ig+\chi_{\{i\}}\,D_ig$, we have $\hatg(A)=\widehat{D_ig}(A\setminus\{i\})$ for
$A\ni i$. Hence $\Inf_i(f)=\E[(D_ig)^2]=\Pr[\mathcal P_i]$, and the law $\nu_i$ of $S\setminus\{i\}$ given $i\in S$
is $\nu_i(T)=\widehat{D_ig}(T)^2/\Inf_i(f)$.

\emph{Gibbs' inequality.} If $Q\ge0$ is a function on subsets of $[n]$ with $\sum_AQ(A)\le1$ and $Q(A)>0$ whenever
$P(A)>0$, then $H(f)\le\sum_AP(A)\log(1/Q(A))$.

\emph{Edge isoperimetry}~\cite{Har64,Ber67,Har76}. A set $B\subseteq\bits^n$ of density $\alpha$ has at least
$\alpha 2^n\log(1/\alpha)$ edges of the cube leaving it. Since $I(f)=2^{1-n}\cdot\#\{\text{edges between }f^{-1}(0)
\text{ and }f^{-1}(1)\}$, we get $I(f)\ge2\alpha\log(1/\alpha)$ with $\alpha=\min(p,1-p)$.

\begin{lemma}\label{lem:empty}
$h(\Var(f))\le 2.8854\,I(f)$ and $H(|S|)\le3.8854\,I(f)$.
\end{lemma}

\begin{proof}
Let $\alpha=\min(p,1-p)\le1/2$ and $\varepsilon=\Var(f)=4\alpha(1-\alpha)\in[2\alpha,4\alpha]$. Using
$(1-\varepsilon)\ln\frac1{1-\varepsilon}\le\varepsilon$,
\[
h(\varepsilon)\le\varepsilon\log\frac1\varepsilon+\varepsilon\log e\le4\alpha\Bigl(\log\frac1\alpha-1\Bigr)+4\alpha\log e
\le2I(f)+1.7708\,\alpha\le2.8854\,I(f),
\]
where we used $I(f)\ge2\alpha\log(1/\alpha)$ and $I(f)\ge\Var(f)\ge2\alpha$. For the second bound, the chain rule gives
$H(|S|)=h(\varepsilon)+\varepsilon\,H(|S|\mid S\ne\emptyset)$, and for a random variable $K$ with values in
$\{1,2,\dots\}$ Gibbs' inequality against the law $2^{-k}$ gives $H(K)\le\E K$. Hence
$\varepsilon\,H(|S|\mid S\neq\emptyset)\le\varepsilon\,\E[|S|\mid S\ne\emptyset]=I(f)$.
\end{proof}

\section{Read-$k$ DNF formulas}\label{sec:readk}

Fix a read-$k$ DNF formula computing $f$. Term $t$ is satisfied exactly when $x$ agrees with a fixed assignment
$a_t\in\bits^{V_t}$ on $V_t$. Put $p_0=\Pr[f=0]$ and $\mu=\sum_t2^{-|V_t|}$.

\begin{lemma}[Path lemma]\label{lem:path}
$p_0\,\mu\le\frac k4\,I(f)$.
\end{lemma}

\begin{proof}
Fix an order on each $V_t$. For a pair $(x,t)$ with $f(x)=0$, change the coordinates of $V_t$ at which $x$ differs
from $a_t$ one at a time, in this order. The function $f$ is $0$ at the start and $1$ at the end, so the path
contains a first edge $\{y,y\oplus e_j\}$ with $f(y)=0$ and $f(y\oplus e_j)=1$, and $j\in V_t$. Given this edge
and $t$, the point $x$ agrees with $y$ outside the set of coordinates of $V_t$ that precede $j$, so at most
$2^{|V_t|-1}$ pairs $(x,t)$ are mapped to the same edge and the same $t$. Since at most $k$ terms contain $j$,
\[
\sum_{x:\,f(x)=0}\ \sum_t2^{-|V_t|}\le\sum_{\text{boundary edges}}\ \sum_{t:\,j\in V_t}2^{|V_t|-1}\,2^{-|V_t|}
\le\frac k2\cdot2^{n-1}I(f).
\]
Dividing by $2^n$ proves the lemma.
\end{proof}

\begin{lemma}[Witness decomposition]\label{lem:witness}
Let $i\in[n]$ with $\Inf_i(f)>0$, and let $t_1,\dots,t_r$ ($r\le k$) be the terms whose variable sets contain $i$.
For $\beta>1$ let $E_i$ be the set of $T\subseteq[n]\setminus\{i\}$ such that
$|(T\cup\{i\})\cap V_{t_s}|<|V_{t_s}|/\beta$ for every $s$. Then
\[
\nu_i(E_i)\le\sum_{s=1}^r\Pr\Bigl[\Bin\bigl(|V_{t_s}|-1,\tfrac12\bigr)<\tfrac{|V_{t_s}|}{\beta}-1\Bigr].
\]
\end{lemma}

\begin{proof}
If $i$ is pivotal at $x$, some term is satisfied at exactly one of $x^{i\leftarrow0},x^{i\leftarrow1}$. That term
depends on $x_i$, so it is some $t_s$, and $x$ lies in the subcube
$C_s=\{x:x_j=a_{t_s}(j)\text{ for }j\in W_s\}$, where $W_s=V_{t_s}\setminus\{i\}$. Let $\mathcal P_{i,s}$ be the
set of pivotal points whose least such index is $s$, and put $d_s=D_ig\cdot1_{\mathcal P_{i,s}}$. The sets
$\mathcal P_{i,s}$ partition $\mathcal P_i$, so $D_ig=\sum_sd_s$ and $\Inf_i(f)=\sum_s\lambda_s\Inf_i(f)$ with
$\lambda_s=\|d_s\|_2^2/\Inf_i(f)$.

Write $x=(y,z)$ with $y$ the restriction of $x$ to $W_s$. Because $\mathcal P_{i,s}\subseteq C_s$,
\[
d_s(y,z)=\mathbf 1[y=a_{t_s}|_{W_s}]\cdot\psi_s(z),\qquad \psi_s(z)=d_s(a_{t_s}|_{W_s},z).
\]
The two factors depend on disjoint sets of coordinates, so Fourier coefficients multiply, and every Fourier
coefficient of the indicator of a point of $\bits^{W_s}$ has absolute value $2^{-|W_s|}$. Therefore
$\widehat{d_s}(T)^2=4^{-|W_s|}\,\widehat{\psi_s}(T\setminus W_s)^2$: under the probability measure
$\nu_{i,s}(T)=\widehat{d_s}(T)^2/\|d_s\|_2^2$ the set $T\cap W_s$ is uniformly distributed over the subsets of $W_s$.
By the Cauchy--Schwarz inequality, for every $T$,
$\bigl(\sum_s\widehat{d_s}(T)\bigr)^2\le\bigl(\sum_s\lambda_s\bigr)\bigl(\sum_s\widehat{d_s}(T)^2/\lambda_s\bigr)$,
the sum running over $s$ with $\lambda_s>0$. Summing over $T\in E_i$ gives $\nu_i(E_i)\le\sum_s\nu_{i,s}(E_i)$.
Finally $E_i\subseteq\{T:|T\cap W_s|<|V_{t_s}|/\beta-1\}$, and $|T\cap W_s|$ has the law $\Bin(|W_s|,1/2)$ under
$\nu_{i,s}$.
\end{proof}

For $\beta>2$ put
\[
C_\beta=\sup_{m\ge1}\,m\,\Pr\Bigl[\Bin\bigl(m-1,\tfrac12\bigr)<\tfrac m\beta-1\Bigr],
\]
which is finite by the Chernoff bound. A direct computation gives $C_3=13\cdot299/4096=0.9489\ldots$, attained at
$m=13$.

\begin{theorem}\label{thm:readk-main}
Let $f$ be computed by a read-$k$ DNF formula and let $\beta>2$. Then
\[
H(f)\le\bigl(2.886+k\,(2+2\beta+2C_\beta)\bigr)\,I(f).
\]
With $\beta=3$ this gives $H(f)\le(2.886+9.898\,k)\,I(f)$.
\end{theorem}

\begin{proof}
\emph{Reference measure.} For a set $U$ of terms let $V(U)$ be the union of their variable sets and
$q(U)=\prod_{t\in U}\pi_t\prod_{t\notin U}(1-\pi_t)$ with $\pi_t=2^{-|V_t|}$. Put $\varepsilon=\Var(f)$, $Q(\emptyset)=1-\varepsilon$ and,
for $A\ne\emptyset$, $Q(A)=\varepsilon\,Q_1(A)$ with
\[
Q_1(A)=\sum_Uq(U)\,2^{-|V(U)|}\,\mathbf 1[A\subseteq V(U)].
\]
Summing first over $A\subseteq V(U)$ shows $\sum_AQ(A)\le1$. A variable that lies in no term is irrelevant to $f$,
so $Q(A)>0$ whenever $P(A)>0$. Gibbs' inequality gives
\begin{equation}\label{eq:gibbs}
H(f)\le h(\varepsilon)+\sum_{A\ne\emptyset}P(A)\log\frac1{Q_1(A)}.
\end{equation}
If $U$ covers $A$, then $Q_1(A)\ge q(U)2^{-|V(U)|}$, so
\begin{equation}\label{eq:cover}
\log\frac1{Q_1(A)}\le\Lambda+\sum_{t\in U}\Bigl(\log\frac1{\pi_t}+|V_t|\Bigr)=\Lambda+2\sum_{t\in U}|V_t|,
\qquad \Lambda=\sum_t\log\frac1{1-\pi_t}\le2\mu ,
\end{equation}
where the last inequality holds because $\pi_t\le\frac12$ and $\log\frac1{1-x}\le2x$ for $0\le x\le\frac12$ by
convexity.

\emph{The cover.} For $A\ne\emptyset$ let $D(A)$ be the set of terms with $|A\cap V_t|\ge|V_t|/\beta$. Call $i\in A$
\emph{sparse} if no term in $D(A)$ contains it, and let $w_i$ be the smallest size of a term containing $i$. Cover
$A$ by $D(A)$ together with a smallest term through each sparse element. Since every variable lies in at most $k$
terms, $\sum_{t\in D(A)}|V_t|\le\beta\sum_{t\in D(A)}|A\cap V_t|\le\beta k|A|$, and \eqref{eq:cover} gives
\[
\log\frac1{Q_1(A)}\le\Lambda+2\beta k|A|+2\sum_{i\in A\text{ sparse}}w_i .
\]
Now $\sum_{A\ni i}P(A)\,\mathbf 1[i\text{ sparse in }A]=\Inf_i(f)\,\nu_i(E_i)$ with $E_i$ as in \cref{lem:witness},
and a sparse $i$ lies only in terms of size larger than $\beta$. By \cref{lem:witness} and $w_i\le|V_{t_s}|$,
\[
w_i\,\nu_i(E_i)\le\sum_s|V_{t_s}|\Pr\Bigl[\Bin\bigl(|V_{t_s}|-1,\tfrac12\bigr)<\tfrac{|V_{t_s}|}\beta-1\Bigr]\le k\,C_\beta .
\]
Summing over $A\ne\emptyset$,
\[
\sum_{A\ne\emptyset}P(A)\log\frac1{Q_1(A)}\le\varepsilon\Lambda+2\beta k\,I(f)+2kC_\beta\,I(f).
\]

\emph{The vacuum term.} By \cref{lem:path}, $\varepsilon\Lambda\le4p_0\cdot2\mu\le2k\,I(f)$.

\emph{The empty set.} $h(\varepsilon)\le2.886\,I(f)$ by \cref{lem:empty}. Adding the three bounds proves the theorem.
\end{proof}

\begin{proof}[Proof of \cref{thm:readk}]
For DNF formulas the first bound is \cref{thm:readk-main} with $\beta=3$, and the second is \cref{thm:readk-log},
proved in \cref{sec:readk-log}. If $f$ is computed by a read-$k$ CNF formula, the dual
function $f^\ast(x)=1-f(\bar x)$, where $\bar x$ is the complement of $x$, is computed by a read-$k$ DNF formula. Its
$\pm1$ version is $x\mapsto-g(\bar x)$, whose Fourier coefficients are $\pm\hatg(A)$, so $f$ and $f^\ast$ have the
same spectral distribution and the same total influence.
\end{proof}

\begin{remark}\label{rem:readk}
(1) Averaging the proof of \cref{lem:path} over uniformly random orders of the sets $V_t$ gives
$p_0\mu\le\frac12\max_j\bigl(\sum_{t\ni j}|V_t|^{-1}\bigr)I(f)$, which improves the vacuum term when all terms are
long. (2) For monotone formulas one can also take $\pi_t$ to be the probability that $t$ is the unique satisfied
term divided by the probability that no other term is satisfied; Harris' inequality then bounds $\log(1/\pi_t)$ by
$|V_t|$ plus the number of terms meeting $V_t$, which leads to a constant of order $k^2$. (3) The dependence on $k$
enters in three places: the vacuum term, the Cauchy--Schwarz step, and the cover. For the star
$x_c\wedge(x_1\vee\cdots\vee x_k)$, written as the read-$k$ formula $\bigvee_j(x_c\wedge x_j)$, one has
$p_0\mu\approx k/8$ while $I(f)\approx1$, so \cref{lem:path} is sharp up to a factor $2$, and the reference measure of
the proof loses a factor of order $k$ on this formula, which is a read-once formula.
\end{remark}

\subsection{Logarithmic dependence on $k$}\label{sec:readk-log}

The reference measure of \cref{thm:readk-main} pays for every term of the cover twice, once to activate it and once
for the variables it adds to $V(U)$, and the dense terms through a single variable can be as many as $k$. The
following reference measure activates only one term in each connected piece of the cover, reaches the other terms of
that piece through shared variables at a logarithmic price, and pays for the subset of $V(U)$ according to its
density.

For $\beta>2$ and integers $w\ge1$ put $\tau_\beta(m)=\Pr[\Bin(m-1,\frac12)<\frac m\beta-1]$,
\[
T_\beta(w)=\sup_{m\ge w}m\,\tau_\beta(m),\qquad M_\beta(k)=\sup_{w\ge1}\min\bigl(w,\,k\,T_\beta(w)\bigr).
\]
By the Chernoff bound $\tau_\beta(m)\le e^{-(m-1)D_\beta}$, where $D_\beta=\frac1\beta\ln\frac2\beta+
(1-\frac1\beta)\ln\bigl(2-\frac2\beta\bigr)$ is the relative entropy of the Bernoulli laws with means $\frac1\beta$ and
$\frac12$. The function $m\mapsto me^{-(m-1)D_\beta}$ increases up to $1/D_\beta$ and decreases afterwards, so
\begin{equation}\label{eq:Mbound}
M_\beta(k)\le\max\Bigl(\frac1{D_\beta},\,1+\frac{\ln k}{D_\beta}\Bigr).
\end{equation}
Indeed, if $w\ge w_0=1+\ln k/D_\beta$ and $w_0\ge1/D_\beta$, then $k\,T_\beta(w)\le kwe^{-(w-1)D_\beta}\le w_0$, while
$\min(w,kT_\beta(w))\le w$ always.

\begin{theorem}\label{thm:readk-log}
Let $f$ be computed by a read-$k$ DNF or CNF formula and let $\beta>2$. Then
\[
H(f)\le\Bigl(11.39+\beta+\log k+\log(\beta k+1)+\log(\beta k)+\kappa_\beta M_\beta(k)\Bigr)\,I(f),
\]
where $\kappa_\beta=1+(\log e)\bigl(\frac1{\beta k}+\frac1{\beta k+1}\bigr)$. In particular
$H(f)\le(30+5.6\log k)\,I(f)$ for every $k\ge1$.
\end{theorem}

\begin{proof}
By duality, as in the proof of \cref{thm:readk}, it suffices to treat DNF formulas. Put $\lambda=1/k$,
$\pi_t=2^{-|V_t|}$ and $\varepsilon=\Var(f)$.

\emph{Reference measure.} Consider the following random procedure, which may fail.
(a) Choose a set $R$ of terms, containing each term $t$ independently with probability $\lambda\pi_t/(1+\lambda\pi_t)$.
(b) Put the terms of $R$ into a queue. When a term $u$ is taken from the queue, choose an integer $c\ge0$ with
probability $1/((c+1)(c+2))$, a multiset of $c$ variables of $V_u$ uniformly among the $\binom{|V_u|+c-1}{c}$ such
multisets, and for each chosen variable $j$ a term containing $j$, uniformly among the at most $k$ such terms; append
the chosen terms to the queue. Let $U$ be the set of terms taken from the queue once it is empty (the procedure fails
if it never empties).
(c) Choose $s\ge1$ with probability $1/(s(s+1))$ and a uniformly random $s$-element subset of $V(U)$ (fail if
$s>|V(U)|$).
Let $Q_1(A)$ be the probability that the procedure outputs the set $A$; then $\sum_AQ_1(A)\le1$. If a nonempty set $A$
is the output of a particular run, $Q_1(A)$ is at least the probability of that run, and so
\begin{multline}\label{eq:animal}
\log\frac1{Q_1(A)}\le\Lambda_\lambda+\sum_{r\in R}\log\frac1{\lambda\pi_r}
+\sum_{u\in U}\Bigl(\log\bigl((c_u+1)(c_u+2)\bigr)+\log\binom{|V_u|+c_u-1}{c_u}+c_u\log k\Bigr)\\
+\log\bigl(s(s+1)\bigr)+\log\binom{|V(U)|}{s},
\end{multline}
where $c_u$ is the integer chosen when $u$ is processed and $\Lambda_\lambda=\sum_t\log(1+\lambda\pi_t)\le(\log e)\,\mu/k$.

\emph{A run producing $A$.} Let $A\ne\emptyset$ with $P(A)>0$, let $D(A)$ be the set of terms with
$|A\cap V_t|\ge|V_t|/\beta$, and call $i\in A$ sparse if no term of $D(A)$ contains it, as in the proof of
\cref{thm:readk-main}. Let $U_d\subseteq D(A)$ be a minimal family covering the non-sparse elements of $A$, and for
each sparse $i$ let $t_i$ be a smallest term containing $i$, of size $w_i$. Put $U=U_d\cup\{t_i\}$. By minimality every
term of $U_d$ contains a non-sparse element of $A$ that no other term of $U_d$ contains, and no $t_i$ belongs to
$D(A)$; hence $|U|\le|A|$. In the graph on $U$ in which two terms are adjacent when their variable sets meet, choose
in each connected component a root and a breadth-first spanning tree, and let $R$ be the set of roots. Processing the
terms along these trees, each non-root term being chosen from its parent through a common variable, and then choosing
$s=|A|$ and the subset $A$ of $V(U)$, is a run that outputs $A$; in it $c_u$ is the number of children of $u$, so
$\sum_uc_u=|U|-|R|$.

\emph{Cost of the run.} Put $s=|A|$ and $W=\sum_{i\in A\text{ sparse}}w_i$. The roots lie in different components,
so their variable sets are pairwise disjoint; a root in $U_d$ has $|V_r|\le\beta|A\cap V_r|$, and every other root is
some $t_i$. Hence
\[
\sum_{r\in R}\log\frac1{\lambda\pi_r}=\sum_{r\in R}\bigl(\log k+|V_r|\bigr)\le|R|\log k+\beta s+W.
\]
Since $(c+1)(c+2)\le2\cdot3^c$, the first sum over $u$ in \eqref{eq:animal} contributes at most $|U|+(|U|-|R|)\log3
\le s(1+\log3)$, and the terms $c_u\log k$ together with $|R|\log k$ give $|U|\log k\le s\log k$. Every variable lies
in at most $k$ terms, so $V_{\mathrm{tot}}=\sum_{u\in U}|V_u|\le\beta\sum_{u\in U_d}|A\cap V_u|+W\le\beta ks+W$. With
$C=\sum_uc_u\le s$, the bound $\log\binom{n+c-1}{c}\le c\log\frac{e(n+c)}c$ and the log-sum inequality give
\[
\sum_{u\in U}\log\binom{|V_u|+c_u-1}{c_u}\le C\log\frac{e(V_{\mathrm{tot}}+C)}{C}\le s\log\frac{e(V_{\mathrm{tot}}+s)}s
\le s\log\bigl(e(\beta k+1)\bigr)+\frac{(\log e)\,W}{\beta k+1},
\]
because $C\mapsto C\log(e(V+C)/C)$ is increasing and $\log(a+x)\le\log a+x\log e/a$. Finally
$\log(s(s+1))\le s\log3$, and $|V(U)|\le V_{\mathrm{tot}}$ gives
$\log\binom{|V(U)|}s\le s\log\frac{e|V(U)|}s\le s\log(e\beta k)+(\log e)W/(\beta k)$. Altogether
\[
\log\frac1{Q_1(A)}\le\Lambda_\lambda+|A|\bigl(\log k+\beta+1+2\log3+2\log e+\log(\beta k+1)+\log(\beta k)\bigr)
+\kappa_\beta W .
\]

\emph{Summation.} As in the proof of \cref{thm:readk-main}, $\sum_{A\ne\emptyset}P(A)W(A)=\sum_i\Inf_i(f)\,w_i\,
\nu_i(E_i)$. By \cref{lem:witness}, $\nu_i(E_i)\le\min(1,k\max_s\tau_\beta(|V_{t_s}|))$, and $w_i\tau_\beta(m)\le
T_\beta(w_i)$ for $m\ge w_i$, so $w_i\nu_i(E_i)\le\min(w_i,kT_\beta(w_i))\le M_\beta(k)$. By \cref{lem:path},
$\varepsilon\Lambda_\lambda\le4p_0(\log e)\mu/k\le(\log e)\,I(f)$, and $h(\varepsilon)\le2.886\,I(f)$ by
\cref{lem:empty}. Inserting these bounds in \eqref{eq:gibbs}, with $Q_1$ in place of the measure used there, and
using $2.886+\log e+1+2\log3+2\log e<11.39$ proves the first claim.

For the second take $\beta=8$, so that $D_8=0.3163\ldots$ and \eqref{eq:Mbound} gives $M_8(1)\le3.17$ and
$M_8(k)\le1+2.2\log k$ for $k\ge2$. Since $\log(8k+1)\le3.17+\log k$ and $\log(8k)=3+\log k$, the constant is at most
$25.56+1.35\cdot3.17<30$ for $k=1$, and at most $25.56+3\log k+1.18\,(1+2.2\log k)<30+5.6\log k$ for $k\ge2$.
\end{proof}

Optimising $\beta$ for each $k$ improves the constant of \cref{thm:readk-log} to less than $20+4.44\log k$ for
$k\le10^6$. The bound of \cref{thm:readk-main} is smaller for $k\le2$. The reference measure does not need the
activation of the dense terms, which is the source of the factor $k$ in \cref{thm:readk-main}: once one term of a
connected piece of the cover is activated, a term overlapping it is described by a variable of an already chosen term
and an index among at most $k$ terms, and the variables it brings are paid for only through the density of $A$ in
$V(U)$, which is at least $1/(\beta k+1)$ up to the sparse terms.

\section{Regular functions}\label{sec:regular}

Negating an input multiplies $\hatg(A)$ by $\pm1$, so we may assume that $f$ is monotone and that its coordinates
are ordered by decreasing desirability: $f(x^{i\leftarrow1,j\leftarrow0})\ge f(x^{i\leftarrow0,j\leftarrow1})$
whenever $i<j$. For $i\in[n]$ let $U_i=\{1,\dots,i-1\}$.

\begin{lemma}\label{lem:antichain}
For every $i$ and every assignment $z$ to the coordinates outside $U_i\cup\{i\}$, the set of $y\in\bits^{U_i}$ such
that $i$ is pivotal at $(y,z)$ is an antichain of $\bits^{U_i}$.
\end{lemma}

\begin{proof}
Suppose $y<y'$ both have $i$ pivotal, choose $j\in U_i$ with $y_j=0<y'_j$, and let $y''=y\oplus e_j\le y'$. By
monotonicity and since $j<i$,
\[
f(y',z,x_i=0)\ge f(y'',z,x_i=0)\ge f(y,z,x_i=1)=1,
\]
so $i$ is not pivotal at $(y',z)$.
\end{proof}

\begin{lemma}[Littlewood--Offord form of Sperner's theorem]\label{lem:sperner}
If $\mathcal A\subseteq\bits^m$ is an antichain, then $|\mathcal A|\,2^{-m}\le\binom m{\lfloor m/2\rfloor}2^{-m}
\le\frac{2}{\sqrt\pi}(m+1)^{-1/2}$ for every $m\ge0$.
\end{lemma}

\begin{proof}
The first inequality is Sperner's theorem~\cite{Spe28}. For even $m=2r\ge2$, $\binom{2r}r4^{-r}\le(\pi r)^{-1/2}
=(2/(\pi m))^{1/2}$. For odd $m=2r+1\ge3$, $\binom{2r+1}r2^{-2r-1}=\binom{2r}r4^{-r}\frac{2r+1}{2r+2}$ and
$(2r+1)^3\le2r(2r+2)^2$, which again gives $\le(2/(\pi m))^{1/2}$; for $m=1$ the bound $1/2\le(2/\pi)^{1/2}$ is
direct. Finally $(2/(\pi m))^{1/2}\le\frac2{\sqrt\pi}(m+1)^{-1/2}$ for $m\ge1$, and for $m=0$ the claim reads
$1\le2/\sqrt\pi$.
\end{proof}

\begin{lemma}[Partial noise sensitivity of the derivative]\label{lem:noise}
Let $i\ge2$, $N=i$, $0<c\le N$ and $\rho=1-c/N$. Then
\[
\sum_{T}\widehat{D_ig}(T)^2\,\rho^{|T\cap U_i|}\le\frac2{\sqrt\pi}\,c^{-1/2}\,\Inf_i(f).
\]
\end{lemma}

\begin{proof}
Write $d=D_ig$. Let $x$ be uniform, let $R\subseteq U_i$ contain each $j\in U_i$ independently with probability
$c/N$, and let $y$ agree with $x$ outside $R$ and be uniform on $R$, independently of $x$. Then
$\E[\chi_T(x)\chi_{T'}(y)]=\mathbf 1[T=T']\rho^{|T\cap U_i|}$, so the left side equals
$\E[d(x)d(y)]\le\E\bigl[|d(x)|\Pr[d(y)\ne0\mid x]\bigr]$. Given $x$ and $R$, the points $y_R$ with $d(y)\ne0$ form
an antichain of $\bits^R$ by \cref{lem:antichain} (the restriction of an antichain to a subcube is an antichain),
so \cref{lem:sperner} gives $\Pr[d(y)\ne0\mid x,R]\le\frac2{\sqrt\pi}(|R|+1)^{-1/2}$. The size $|R|$ has the law
$\Bin(N-1,c/N)$, so $\E[1/(|R|+1)]=(1-(1-c/N)^N)/c\le1/c$ and, by Jensen's inequality,
$\E[(|R|+1)^{-1/2}]\le c^{-1/2}$. Finally $\E|d(x)|=\Inf_i(f)$.
\end{proof}

\begin{lemma}\label{lem:ranktail}
Let $B=3\sqrt{3/\pi}$ and $L_1=2\ln B+2=4.1511\ldots$. For every $i$ with $\Inf_i(f)>0$,
\[
\E_{T\sim\nu_i}\Bigl[\ln\frac{i}{|T\cap U_i|+1}\Bigr]\le L_1 .
\]
\end{lemma}

\begin{proof}
For $i=1$ the left side is $0$. Let $i\ge2$, $N=i$ and $\kappa(T)=|T\cap U_i|\in[0,N-1]$. For real $K\ge1$ apply
\cref{lem:noise} with $c=N/(3K)$, so that $\rho=1-1/(3K)$:
\[
\nu_i(\kappa\le K)\le\rho^{-K}\,\E_{\nu_i}\rho^{\kappa}\le\Bigl(1-\frac1{3K}\Bigr)^{-K}\frac{2}{\sqrt\pi}
\Bigl(\frac{3K}N\Bigr)^{1/2}\le B\Bigl(\frac KN\Bigr)^{1/2},
\]
because $-K\ln(1-\frac1{3K})=\sum_{m\ge1}3^{-m}K^{1-m}/m$ is nonincreasing in $K$ and equals $\ln\frac32$ at $K=1$.
Since $\ln(N/(\kappa+1))\in[0,\ln N]$,
\[
\E_{\nu_i}\ln\frac N{\kappa+1}=\int_0^{\ln N}\nu_i\bigl(\kappa+1<Ne^{-u}\bigr)\,du
\le\int_0^\infty\min\bigl(1,Be^{-u/2}\bigr)\,du=2\ln B+2,
\]
using the previous bound with $K=Ne^{-u}\ge1$.
\end{proof}

\begin{lemma}\label{lem:ranksum}
$\E_S\bigl[\sum_{i\in S}\ln^+(i/|S|)\bigr]\le L_1\,I(f)$.
\end{lemma}

\begin{proof}
For $i\in S$ and $T=S\setminus\{i\}$ we have $|T\cap U_i|+1\le\min(|S|,i)$, hence
$\ln^+(i/|S|)\le\ln(i/(|T\cap U_i|+1))$. Since $\Pr[S=T\cup\{i\}]=\Inf_i(f)\nu_i(T)$, summing \cref{lem:ranktail}
with weights $\Inf_i(f)$ proves the claim.
\end{proof}

\begin{lemma}[Rank-weighted reference measure]\label{lem:rankref}
For every random subset $S$ of $[n]$ and every $a>1$,
\[
H(S\mid|S|)\le\frac1{\ln2}\Bigl(\ln\frac{ea}{a-1}\,\E|S|+a\,\E\sum_{i\in S}\ln^+\frac i{|S|}\Bigr).
\]
\end{lemma}

\begin{proof}
For $k\ge1$ let $\theta_k(s)=\min(1,(k/s)^a)$ and let $\pi_k(A)=Z_k^{-1}\prod_{i\in A}\theta_k(i)$ on $k$-sets,
where $Z_k$ is the $k$-th elementary symmetric polynomial of $\theta_k(1),\dots,\theta_k(n)$. Then
$Z_k\le(\sum_s\theta_k(s))^k/k!$, $\sum_s\theta_k(s)\le k+\int_k^\infty(k/s)^a\,ds=\frac{ka}{a-1}$ and
$k!\ge(k/e)^k$, so $\ln Z_k\le k\ln\frac{ea}{a-1}$. Gibbs' inequality on each level gives
$H(S\mid|S|=k)\le\frac1{\ln2}\E\bigl[\ln Z_k+a\sum_{i\in S}\ln^+(i/k)\bigm||S|=k\bigr]$; average over $k$.
\end{proof}

\begin{proof}[Proof of \cref{thm:regular}]
By the chain rule, \cref{lem:empty}, \cref{lem:rankref} and \cref{lem:ranksum},
\[
H(f)=H(|S|)+H(S\mid|S|)\le\Bigl(3.8854+\frac{1+\ln\frac a{a-1}+aL_1}{\ln2}\Bigr)I(f).
\]
The bracket is minimised at $a=\frac12\bigl(1+\sqrt{1+4/L_1}\bigr)=1.2006\ldots$, where the second summand equals
$11.214\ldots$, giving $H(f)\le15.1\,I(f)$.
\end{proof}

The same argument applies to Boolean functions of several regular functions, even when their desirability orders
are unrelated.

\begin{theorem}\label{thm:multi}
Let $h_1,\dots,h_m$ be regular Boolean functions on $\bits^n$, with arbitrary and possibly different desirability
orders and orientations, and let $F:\bits^m\to\bits$ be arbitrary. Then $f=F(h_1,\dots,h_m)$ satisfies
\[
H(f)\le(16.5+8.4\ln m)\,I(f).
\]
In particular this holds for every intersection of $m$ halfspaces, and for $x\mapsto G(w\cdot x)$ when
$G:\mathbb R\to\bits$ changes sign at most $m$ times.
\end{theorem}

\begin{proof}
For each $l$ fix an orientation and a desirability order of $h_l$, let $r_l(i)$ be the position of $i$ in the
$l$-th order and let $U_i^l$ be the set of the $r_l(i)-1$ coordinates that precede $i$ in it. If $i$ is pivotal for
$f$ at $x$, then it is pivotal for some $h_l$ at $x$. Let $d_{i,l}=D_ig\cdot\mathbf 1[\,l$ is the least index for
which $i$ is pivotal for $h_l\,]$; these functions have disjoint supports, $D_ig=\sum_ld_{i,l}$, and the support of
$d_{i,l}$ lies in the pivotal set of $i$ for $h_l$. \Cref{lem:antichain} applied to $h_l$ shows that this support,
restricted to any subcube spanned by coordinates of $U_i^l$, is an antichain in that subcube. The proof of
\cref{lem:noise} uses nothing else, so it applies to $d_{i,l}$ with $U_i^l$ in place of $U_i$, and the proof of
\cref{lem:ranktail} then gives, for the probability measure $\nu_{i,l}(T)=\widehat{d_{i,l}}(T)^2/\|d_{i,l}\|_2^2$ and
$F_l(T)=\ln\frac{r_l(i)}{|T\cap U_i^l|+1}$,
\[
\nu_{i,l}(F_l>u)\le\min(1,Be^{-u/2})\quad(u\ge0),\qquad\text{hence}\qquad
\E_{\nu_{i,l}}e^{\lambda F_l}\le\frac{B^{2\lambda}}{1-2\lambda}\quad(0<\lambda<\tfrac12).
\]
As in the proof of \cref{lem:witness}, the Cauchy--Schwarz inequality gives $\nu_i\le\sum_l\nu_{i,l}$. For
$G_i=\min_lF_l$, Jensen's inequality and this domination give, with $\lambda=\frac14$,
\[
\E_{\nu_i}G_i\le\frac1\lambda\ln\E_{\nu_i}e^{\lambda G_i}\le\frac1\lambda\ln\sum_l\E_{\nu_{i,l}}e^{\lambda F_l}
\le4\ln m+4\ln2+2\ln B<4\ln m+4.93 .
\]
Since $\ln^+\frac{r_l(i)}{|S|}\le F_l(S\setminus\{i\})$ for $i\in S$, summing with weights $\Inf_i(f)$ as in the
proof of \cref{lem:ranksum} gives $\E_S\sum_{i\in S}\min_l\ln^+\frac{r_l(i)}{|S|}\le(4\ln m+4.93)\,I(f)$.

Now repeat the proof of \cref{lem:rankref} with $\theta_k(i)=\max_l\min\bigl(1,(k/r_l(i))^a\bigr)$. Each order is a
permutation of $[n]$, so $\sum_i\theta_k(i)\le mka/(a-1)$ and $\ln Z_k\le k\bigl(1+\ln m+\ln\frac a{a-1}\bigr)$, while
$\ln\frac1{\theta_k(i)}=a\min_l\ln^+\frac{r_l(i)}k$. Together with \cref{lem:empty},
\[
H(f)\le\Bigl(3.886+\frac{1+\ln\frac a{a-1}+4.93\,a+(1+4a)\ln m}{\ln2}\Bigr)I(f),
\]
and $a=1.2$ gives the stated bound. For $G(w\cdot x)$, write $f$ as a Boolean function of the $m$ threshold
functions $\mathrm{sgn}(w\cdot x-\theta_l)$ at the sign changes of $G$.
\end{proof}

When the orders of $h_1,\dots,h_m$ coincide, the supports of the pieces lie in a union of $m$ antichains of the same
predecessor cube, the bound of \cref{lem:sperner} is simply multiplied by $m$, and the constant improves to
$3.8854+\min_{a>1}\bigl(1+\ln\frac a{a-1}+a(L_1+2\ln m)\bigr)/\ln2=15.1+O(\log m)$.

\begin{remark}\label{rem:multi-examples}
\Cref{thm:multi} applies to decision lists whose tests are $m$ threshold functions, to unions and intersections of
$m$ halfspaces, and to DNF formulas with $m$ terms, each term being a threshold function; for read-$k$ DNF formulas
with few terms it can be stronger than \cref{thm:readk}. It also applies to every function $f$ that becomes regular
after fixing the values of a set $K$ of $k$ coordinates (the order may depend on the values): $f$ is a Boolean
function of the $k$ coordinates in $K$ and of the $2^k$ restrictions $f(a,\cdot)$, $a\in\bits^K$, all of which are
regular, so $H(f)\le\bigl(16.5+8.4\ln(k+2^k)\bigr)I(f)$.
\end{remark}

\begin{remark}\label{rem:width}
The proof of \cref{thm:regular} uses the order only through \cref{lem:antichain}, which needs every coordinate of
$U_i$ to be at least as desirable as $i$, and through the bound $\sum_s\theta_k(s)\le ka/(a-1)$ in
\cref{lem:rankref}. Let $f$ be monotone, let $\preceq$ be its desirability preorder, refined inside each class of
equivalent coordinates by a fixed order, and let $w$ be the largest number of pairwise incomparable coordinates. Take
$U_i$ to be the set of coordinates above $i$ and $r(i)=|U_i|+1$. \Cref{lem:antichain,lem:noise,lem:ranktail,lem:ranksum}
hold with the same proofs. By Dilworth's theorem the coordinates form $w$ chains, and the $p$-th element of a chain
has $r\ge p$, so $\sum_i\min(1,(k/r(i))^a)\le wka/(a-1)$ and $\ln Z_k\le k\bigl(\ln\frac{ea}{a-1}+\ln w\bigr)$.
Hence $H(f)\le(15.1+\log w)\,I(f)$. A monotone function invariant under $S_{n_1}\times\cdots\times S_{n_w}$ acting on
$w$ blocks of coordinates has width at most $w$; for these functions, with or without monotonicity, O'Donnell, Wright
and Zhou~\cite{OWZ11} proved the better bound $H(f)\le(12.04+\log w)\,I(f)$.
\end{remark}

\begin{remark}
O'Donnell, Wright and Zhou~\cite[Theorem~5]{OWZ11} proved $H(|S|)\le3I(f)$ for every Boolean function. Using it in
place of the second bound of \cref{lem:empty}, the constant of \cref{thm:regular} becomes $14.22$, and those of
\cref{thm:multi} and \cref{rem:width} decrease accordingly. The function
$\mathrm{AND}_n$ is a threshold function with $H/I\to4$, so the optimal constant for regular functions is at least
$4$.
\end{remark}

\section{Unions of regular functions}\label{sec:union}

The two mechanisms combine. In this section $f=h_1\vee\cdots\vee h_m$, where $h_l$ depends exactly on the variables
in a set $R_l$, and every variable lies in at most $k$ of the sets $R_l$.

\begin{theorem}\label{thm:union}
If $h_1,\dots,h_m$ are regular and every variable is relevant to at most $k$ of them, then $h_1\vee\cdots\vee h_m$
and $h_1\wedge\cdots\wedge h_m$ satisfy $H(f)\le(3.9+21.2k)\,I(f)$.
\end{theorem}

The terms of a DNF formula are threshold functions, so \cref{thm:union} contains \cref{thm:readk,thm:regular} up to
the constants. It also applies to unions and intersections of halfspaces in which every variable has a nonzero
weight in at most $k$ of the halfspaces, with a constant that does not depend on their number.

\begin{lemma}\label{lem:upath}
Let $h_1,\dots,h_m$ be arbitrary Boolean functions, $h_l$ depending only on the variables in $R_l$, with every
variable in at most $k$ of the sets $R_l$. For $f=h_1\vee\cdots\vee h_m$ and $p_0=\Pr[f=0]$,
\[
p_0\sum_l\Pr[h_l=1]\le\frac k4\,I(f).
\]
\end{lemma}

\begin{proof}
Fix $l$, let $x$ be uniform and let $x'$ be obtained from $x$ by resampling the coordinates in $R_l$. Since
$h_l(x')$ depends only on $x'_{R_l}$, which is independent of $x$, and $f\ge h_l$,
$\Pr[f(x)=0,\,f(x')=1]\ge p_0\Pr[h_l=1]$. Now resample the coordinates of $R_l$ one at a time. Every intermediate
point is uniform, and resampling a coordinate $j$ of a uniform point changes $f$ from $0$ to $1$ with probability
$\Inf_j(f)/4$. If $f(x)=0$ and $f(x')=1$, one of these steps does so, so the probability above is at most
$\frac14\sum_{j\in R_l}\Inf_j(f)$. Summing over $l$ proves the lemma.
\end{proof}

In \cref{lem:witness} the subcube structure of a term made the spectral sample fill the witnessing term uniformly.
For a regular function we use instead that its pivotal set is small when the function is rarely true.

\begin{lemma}\label{lem:level}
Let $h$ be a Boolean function depending only on the variables in a set $R\ni j$, let $\alpha=\Inf_j(h)$,
and let $d:\bits^{[n]\setminus\{j\}}\to\{-1,0,1\}$ vanish at every point at which $j$ is not pivotal for $h$. For the
probability measure $\nu(T)=\widehat d(T)^2/\|d\|_2^2$ and every integer $0\le\ell\le\ln(1/\alpha)$,
\[
\nu\bigl(|T\cap R|\le\ell\bigr)\le\varphi_\ell(\alpha),\qquad\text{where }\varphi_0(\alpha)=\alpha\text{ and }
\varphi_\ell(\alpha)=\alpha\Bigl(\frac{2e\ln(1/\alpha)}{\ell}\Bigr)^{\ell}\text{ for }\ell\ge1.
\]
\end{lemma}

\begin{proof}
Write the argument of $d$ as $(y,z)$ with $y$ the restriction to $R\setminus\{j\}$, put $d_z=d(\cdot,z)$ and
$W^{\le\ell}[\phi]=\sum_{|U|\le\ell}\widehat\phi(U)^2$. Parseval's identity in $z$ gives
$\sum_{T:\,|T\cap R|\le\ell}\widehat d(T)^2=\E_zW^{\le\ell}[d_z]$ and $\|d\|_2^2=\E_z\|d_z\|_2^2$. Each $d_z$ takes
values in $\{-1,0,1\}$ and vanishes outside the set of $y$ at which $j$ is pivotal for $h$, a set of measure
$\alpha$, so $\|d_z\|_2^2=\alpha_z\le\alpha$. For $\ell=0$, $W^{\le0}[d_z]=(\E d_z)^2\le\alpha_z^2$. For $\ell\ge1$ with
$\alpha_z>0$ put $\delta=\ell/(2\ln(1/\alpha_z))\le\frac12$. By H\"older's inequality and the hypercontractive
inequality $\|\phi\|_{1+1/\delta}\le\delta^{-\ell/2}\|\phi\|_2$ for $\phi$ of degree at most $\ell$,
\[
W^{\le\ell}[d_z]=\langle d_z,d_z^{\le\ell}\rangle\le\alpha_z^{1/(1+\delta)}\,\delta^{-\ell/2}\,W^{\le\ell}[d_z]^{1/2},
\]
hence $W^{\le\ell}[d_z]\le\alpha_z^{2-2\delta}\delta^{-\ell}=\alpha_z^2\bigl(2e\ln(1/\alpha_z)/\ell\bigr)^\ell$. In both
cases $W^{\le\ell}[d_z]\le\alpha_z\varphi_\ell(\alpha_z)$. The function $\varphi_\ell$ is nondecreasing on
$(0,e^{-\ell}]$, which contains every $\alpha_z$, so $\E_zW^{\le\ell}[d_z]\le\varphi_\ell(\alpha)\,\E_z\alpha_z$.
\end{proof}

\begin{proof}[Proof of \cref{thm:union}]
The complement of $h_1\wedge\cdots\wedge h_m$ is the union of the complements of the $h_l$, which are regular, and a
function and its complement have the same spectral distribution; so it suffices to treat unions. For each $l$ fix
an orientation under which $h_l$ is monotone and an order of $R_l$ by decreasing desirability for $h_l$; let
$r_l(j)$ be the position of $j\in R_l$ and $U^l_j$ the set of the $r_l(j)-1$ variables before it. Put
$\pi_l=\min(\Pr[h_l=1],\frac12)$, $\varepsilon=\Var(f)$, $a=1.44$, $\beta=5$, $c_1=1+\frac a{a-1}$ and
$c_2=1+\frac{c_1}\beta$.

\emph{Reference measure.} For $\kappa\ge1$ let $B^l_\kappa$ be the law of the random subset of $R_l$ that contains
each $j$ independently with probability $\theta^l_\kappa(j)=\frac12\min\bigl(1,(\kappa/r_l(j))^a\bigr)$. Let $L$
contain each index $l$ independently with probability $\pi_l$, for each $l\in L$ choose $\kappa_l\ge1$ with
probability $2^{-\kappa_l}$ and then $A_l\sim B^l_{\kappa_l}$, all independently, and let $Q_1(A)$ be the
probability that $\bigcup_{l\in L}A_l=A$. With $Q(\emptyset)=1-\varepsilon$ and $Q(A)=\varepsilon Q_1(A)$ for
$A\ne\emptyset$ we have $\sum_AQ(A)\le1$, and $Q(A)>0$ whenever $P(A)>0$, so Gibbs' inequality gives
\[
H(f)\le h(\varepsilon)+\sum_{A\ne\emptyset}P(A)\log\frac1{Q_1(A)}.
\]

\emph{A choice for each $A$.} Fix $A\ne\emptyset$ and put $\kappa_l=|A\cap R_l|$ and
$\lambda_l=\frac1\beta\log\frac1{\pi_l}$; call $l$ dense if $\kappa_l\ge\lambda_l$. For $j\in A$ and each $l$ with
$j\in R_l$ put $F^j_l=\ln\frac{r_l(j)}{|A\cap U^l_j|+1}\ge0$ and
\[
G_j(l)=\frac a{\ln2}F^j_l+c_2\log\frac1{\pi_l}\,\mathbf 1[\,l\text{ is not dense}\,],
\]
let $\tau(j)$ minimise $G_j(l)$ over these $l$, and write $G_j=G_j(\tau(j))$. Let $L_A$ consist of the dense indices
and the $\tau(j)$, $j\in A$, and for $l\in L_A$ let $A_l=\{j\in A:\tau(j)=l\}$. Then $A=\bigcup_{l\in L_A}A_l$ and
$\kappa_l\ge1$ for $l\in L_A$. Using $\log\frac1{1-x}\le2x$ for $0\le x\le\frac12$ and
$\sum_{j\in R_l}\theta^l_\kappa(j)\le\frac12\bigl(\kappa+\int_\kappa^\infty(\kappa/r)^a\,dr\bigr)=\frac{\kappa a}{2(a-1)}$,
\[
\log\frac1{B^l_\kappa(A_l)}\le\sum_{j\in A_l}\Bigl(1+a\log^+\frac{r_l(j)}\kappa\Bigr)+\frac{\kappa a}{a-1},
\]
and for $j\in A_l$ we have $\kappa_l\ge|A\cap U^l_j|+1$, so that $\log^+(r_l(j)/\kappa_l)\le F^j_l/\ln2$. With
$\Lambda=\sum_l\log\frac1{1-\pi_l}$ this gives
\[
\log\frac1{Q_1(A)}\le\Lambda+\sum_{l\in L_A}\Bigl(\log\frac1{\pi_l}+c_1\kappa_l\Bigr)
+\sum_{j\in A}\Bigl(1+\frac a{\ln2}F^j_{\tau(j)}\Bigr).
\]
For dense $l$, $\log\frac1{\pi_l}\le\beta\kappa_l$. A nondense $l\in L_A$ equals $\tau(j)$ for some $j\in A$ and
satisfies $\log\frac1{\pi_l}+c_1\kappa_l<c_2\log\frac1{\pi_l}$. As $\sum_l\kappa_l\le k|A|$,
\[
\log\frac1{Q_1(A)}\le\Lambda+(\beta+c_1)\,k|A|+\sum_{j\in A}(1+G_j).
\]

\emph{The pieces.} For $j$ with $\Inf_j(f)>0$ let $l_1,\dots,l_r$ ($r\le k$) be the indices with $j\in R_{l_s}$. If
$j$ is pivotal for $f$ at $x$, then it is pivotal for some $h_{l_s}$ at $x$. As in the proof of \cref{thm:multi},
split $D_jg$ according to the least such $s$ into pieces $d_s$ with disjoint supports, $d_s$ vanishing where $j$ is
not pivotal for $h_{l_s}$, and let $\nu_{j,s}$ be the normalised spectral measure of $d_s$. The Cauchy--Schwarz
inequality gives $\nu_j\le\sum_s\nu_{j,s}$, and $G_j\le G_j(l_s)$ for every $s$, so with $A=T\cup\{j\}$,
\[
\E_{T\sim\nu_j}G_j\le\sum_{s=1}^r\Bigl(\frac a{\ln2}\E_{\nu_{j,s}}F^j_{l_s}+c_2\log\frac1{\pi_{l_s}}\;
\nu_{j,s}\bigl(l_s\text{ is not dense}\bigr)\Bigr).
\]
The proof of \cref{thm:multi} gives $\E_{\nu_{j,s}}F^j_{l_s}\le L_1$. Put $\lambda=\lambda_{l_s}$. If $\lambda\le1$, then
$l_s$ is dense because $j\in A\cap R_{l_s}$. Otherwise $\pi_{l_s}=2^{-\beta\lambda}<\frac12$, and $l_s$ is not dense
exactly when $|T\cap R_{l_s}|\le\ell:=\lceil\lambda\rceil-2$. At a point where $j$ is pivotal for $h_{l_s}$, one of
the two neighbours has $h_{l_s}=1$, so $\Inf_j(h_{l_s})\le2\Pr[h_{l_s}=1]=2^{1-\beta\lambda}$; as
$\ell\le(\beta\lambda-1)\ln2$, \cref{lem:level} gives
\[
\log\frac1{\pi_{l_s}}\;\nu_{j,s}\bigl(l_s\text{ is not dense}\bigr)\le\beta\lambda\,\varphi_\ell\bigl(2^{1-\beta\lambda}\bigr)
\le K:=\max_{\ell\ge0}\,5(\ell+1)\,\varphi_\ell\bigl(2^{-5\ell-4}\bigr)=0.6624\ldots;
\]
here we used that for fixed $\ell$ the middle expression decreases in $\lambda\in(\ell+1,\ell+2]$, since its
logarithmic derivative is $\frac1\lambda-\beta\ln2\,(1-\ell/((\beta\lambda-1)\ln2))<0$, and the maximum is attained at
$\ell=1$.

\emph{Conclusion.} Since $P(A)=\Inf_j(f)\,\nu_j(A\setminus\{j\})$ for $j\in A$, summing over $A$ gives
\[
H(f)\le h(\varepsilon)+\varepsilon\Lambda+(\beta+c_1)k\,I(f)+I(f)+k\Bigl(\frac{aL_1}{\ln2}+c_2K\Bigr)I(f).
\]
By \cref{lem:empty}, $h(\varepsilon)\le2.8854\,I(f)$, and by \cref{lem:upath},
$\varepsilon\Lambda\le4p_0\cdot2\sum_l\pi_l\le2k\,I(f)$. The coefficient of $k$ is
$2+5+4.273+8.624+1.855\cdot0.6624<21.2$.
\end{proof}

The device of \cref{sec:readk-log} also applies here, with the witnessing terms replaced by the functions $h_l$ and
the uniform subsets of $V(U)$ replaced by rank laws.

\begin{theorem}\label{thm:union-log}
Under the hypotheses of \cref{thm:union}, $H(f)\le(54+12.8\log k)\,I(f)$.
\end{theorem}

\begin{proof}
Keep the notation of the proof of \cref{thm:union}, but put $a=1.44$, $\beta=10$ and $\lambda=1/k$, and let
$\theta^l_s(i)=\min(1,(s/r_l(i))^a)$ for $i\in R_l$ and $Z^l_s$ the $s$-th elementary symmetric function of the numbers
$\theta^l_s(i)$, $i\in R_l$; as in the proof of \cref{lem:rankref}, $\log Z^l_s\le s\log\frac{ea}{a-1}$.

\emph{Reference measure.} (a) Choose a set $R$ of indices, containing each $l$ independently with probability
$\lambda\pi_l/(1+\lambda\pi_l)$. (b) Put the indices of $R$ into a queue. When $l$ is taken from the queue, choose
$s\ge1$ with probability $1/(s(s+1))$ and an $s$-element subset $B_l\subseteq R_l$ with probability
$\prod_{i\in B_l}\theta^l_s(i)/Z^l_s$; then choose $c\ge0$ with probability $1/((c+1)(c+2))$, a multiset of $c$
elements of $B_l$ uniformly, and for each chosen element $j$ an index $l'$ with $j\in R_{l'}$, uniformly among the at
most $k$ such indices; append these indices to the queue. (c) When the queue is empty, output $\bigcup_lB_l$. Let
$Q_1(A)$ be the probability that the output is $A$.

\emph{A run producing $A$.} Let $A\ne\emptyset$ with $P(A)>0$. Define $F^j_l$, the dense indices and $G_j(l)$ as in the
proof of \cref{thm:union}, with $c_2$ replaced by $1$, let $\tau(j)$ minimise $G_j(l)$ over the $l$ with $j\in R_l$,
put $A_l=\{j\in A:\tau(j)=l\}$ and let $\mathcal F$ be the set of $l$ with $A_l\ne\emptyset$. Consider the directed
graph on $\mathcal F$ with an edge $l\to l'$ when $l\ne l'$ and $A_l\cap R_{l'}\ne\emptyset$, choose one index in each
strongly connected component that receives no edge from another component, and let $R$ be the set of these roots.
Every index of $\mathcal F$ is reached from a root; take a breadth-first spanning forest along the edges, attach each
non-root $l'$ to its parent $l$ through an element of $A_l\cap R_{l'}$, and let $B_l=A_l$ for $l\in\mathcal F$. This is
a run that outputs $A$, with $|R|+\sum_lc_l=|\mathcal F|\le|A|$, where $c_l$ is the number of children of $l$.

The sets $A\cap R_r$, $r\in R$, are pairwise disjoint. Indeed, if $j\in A\cap R_r\cap R_{r'}$ with $r\ne r'$ and
$l=\tau(j)$, then either $l\notin\{r,r'\}$ and both $l\to r$ and $l\to r'$ are edges, or $l=r$ and $r\to r'$ is an
edge, or $l=r'$ and $r'\to r$ is an edge; in each case two roots lie in one strong component or a root's component
receives an edge from another one, which the choice of roots excludes.

\emph{Cost of the run.} Let $\kappa_l=|A\cap R_l|$ and $N_j=\log\frac1{\pi_{\tau(j)}}\mathbf1[\tau(j)\text{ is not
dense}]$, so that $G_j:=G_j(\tau(j))=\frac a{\ln2}F^j_{\tau(j)}+N_j$. A dense root has $\log\frac1{\pi_r}\le\beta\kappa_r$,
and a root that is not dense is $\tau(j)$ for some $j$; by the disjointness just proved,
$\sum_{r\in R}\log\frac1{\lambda\pi_r}\le|R|\log k+\beta|A|+\sum_{j\in A}N_j$. For $i\in A_l$ with $s=|A_l|$, write
$u=|A\cap U^l_i|$; then $u+1\le\kappa_l$ and $\ln^+(r_l(i)/s)\le F^i_l+\ln(\kappa_l/s)$, so by the log-sum inequality and
$\sum_l\kappa_l\le k|A|$,
\[
\sum_{l\in\mathcal F}\sum_{i\in A_l}\log\frac1{\theta^l_{|A_l|}(i)}\le\frac a{\ln2}\sum_{j\in A}F^j_{\tau(j)}
+a\sum_l|A_l|\log\frac{\kappa_l}{|A_l|}\le\frac a{\ln2}\sum_{j\in A}F^j_{\tau(j)}+a|A|\log k .
\]
The choices of $s$, of $c$, of the multisets and of the indices of the children cost at most $|A|\log3$,
$|A|(1+\log3)$, $|A|\log(2e)$ and $(|\mathcal F|-|R|)\log k$ in total, as in the proof of \cref{thm:readk-log}, and
together with $|R|\log k$ from the roots the indices cost $|\mathcal F|\log k\le|A|\log k$. Hence
\[
\log\frac1{Q_1(A)}\le\Lambda_\lambda+|A|\Bigl((1+a)\log k+\beta+1+2\log3+\log\frac{ea}{a-1}+\log(2e)\Bigr)
+\sum_{j\in A}G_j,
\]
with $\Lambda_\lambda=\sum_l\log(1+\lambda\pi_l)\le(\log e)\sum_l\pi_l/k$.

\emph{The pieces.} Let $d_s$, $\nu_{j,s}$ and $l_s$ be as in the proof of \cref{thm:union}, so that
$\nu_j\le\sum_s\nu_{j,s}$ and $G_j\le G_j(l_s)$ for every $s$. Put $b_s=\log\frac1{\pi_{l_s}}$. The proof of
\cref{thm:multi} gives $\nu_{j,s}(F^j_{l_s}>v)\le Be^{-v/2}$, and the argument of the proof of \cref{thm:union} gives
$\nu_{j,s}(l_s\text{ is not dense})\le g(b_s)$, where $g(b)=\sup\{\varphi_{\ell(b')}(2^{1-b'}):b'\ge b,\ b'>\beta\}$ and
$\ell(b')=\lceil b'/\beta\rceil-2$. On each interval $(10(\ell+1),10(\ell+2)]$ the logarithm of
$2^{0.44b}\varphi_\ell(2^{1-b})$ decreases, and its limits at the left end points are $-4.6,-4.04,-4.26,-4.56,\dots$ for
$\ell=0,1,2,3,\dots$, decreasing from $\ell=1$ on, so $g(b)\le C_g2^{-\gamma b}$ with $\gamma=0.44$ and
$C_g=2^{-4.04}$. Since $\{G_j(l_s)>u\}\subseteq\{\frac a{\ln2}F^j_{l_s}>\frac u2\}\cup\{l_s\text{ not dense},\,b_s>\frac u2\}$,
\[
\nu_j(G_j>u)\le\min\bigl(1,\,kBe^{-u\ln2/(4a)}+k\,C_g2^{-\gamma u/2}\bigr),
\]
and integrating, with $\min(1,x+y)\le\min(1,2x)+\min(1,2y)$ and $2C_g<1$,
\[
\E_{\nu_j}G_j\le\frac{4a}{\ln2}\bigl(\ln(2kB)+1\bigr)+\frac2\gamma\bigl(\log k+\log e\bigr).
\]

\emph{Conclusion.} Sum over $A$ as before. We have $h(\varepsilon)\le2.886\,I(f)$, and
$\varepsilon\Lambda_\lambda\le(\log e)\,I(f)$ because $\varepsilon\le4p_0$ and \cref{lem:upath} gives
$4p_0\sum_l\pi_l\le k\,I(f)$. The constant is at most $53.7+(1+a+4a+2/\gamma)\log k<54+12.8\log k$.
\end{proof}

The bound of \cref{thm:union} is smaller for $k\le3$. Optimising $a$ and $\beta$ for each $k$ and integrating the tail
bounds exactly gives a constant below $61$ for $k\le10$ and below $108$ for $k\le1000$.

The two mechanisms also combine with \cref{thm:multi}: the functions $h_l$ may themselves be Boolean functions of
several regular functions.

\begin{theorem}\label{thm:union-multi}
Let $f=h_1\vee\cdots\vee h_m$ or $f=h_1\wedge\cdots\wedge h_m$, where every $h_l$ is a Boolean function of at most $M$
regular functions, whose desirability orders are arbitrary, and every variable is relevant to at most $k$ of the
functions $h_l$. Then
\[
H(f)\le(54+12.8\log k+6.8\log M)\,I(f)\qquad\text{and}\qquad H^+(f)\le(93+23.1\log k+12.6\log M)\,J(f).
\]
\end{theorem}

For instance, unions and intersections of polytopes with at most $M$ facets in which every variable has a nonzero
weight in at most $k$ of the polytopes satisfy the conjecture with a constant that does not depend on their number,
and so do read-$k$ DNF formulas in which every term is replaced by a Boolean function of at most $M$ threshold
functions on its variables. A Boolean function of $w$ variables is a Boolean function of the $w$ dictators $x_j$,
which are regular, so a union $h_1\vee\cdots\vee h_m$ in which every $h_l$ depends on at most $w$ variables and every
variable is relevant to at most $k$ of the $h_l$ satisfies $H(f)\le(54+12.8\log k+6.8\log w)\,I(f)$, whatever the
functions $h_l$ are.

\begin{proof}
A complement of a Boolean function of $M$ regular functions is again one, so it suffices to treat unions. Write
$h_l=F_l(h_{l,1},\dots,h_{l,M})$ with $h_{l,r}$ regular, depending on the variables of a set $R_{l,r}\subseteq R_l$,
and let $r_{l,r}(i)$ be the position of $i$ in a desirability order of $h_{l,r}$ and $U^{l,r}_i$ the set of variables
before it. The proof of \cref{thm:union-log} applies with three changes.

(i) In the reference measure put $\theta^l_s(i)=\max_r\min\bigl(1,(s/r_{l,r}(i))^a\bigr)$, the maximum over the $r$
with $i\in R_{l,r}$. Then $\sum_i\theta^l_s(i)\le Msa/(a-1)$, and as in the proof of \cref{lem:rankref},
$\log Z^l_s\le s\log\frac{eMa}{a-1}$.

(ii) In the run put $F^j_l=\min_r\ln\bigl(r_{l,r}(j)/(|A\cap U^{l,r}_j|+1)\bigr)$, the minimum over the $r$ with
$j\in R_{l,r}$. For $i\in A_l$ and the $r$ attaining this minimum, $|A\cap U^{l,r}_i|+1\le\kappa_l$, so
$\ln^+(r_{l,r}(i)/|A_l|)\le F^i_l+\ln(\kappa_l/|A_l|)$ and the rank costs are bounded as before. The per-set bound
gains the term $|A|\log M$ from (i).

(iii) If $j$ is pivotal for $h_{l_s}$, it is pivotal for some $h_{l_s,r}$. Splitting the piece $d_s$ further
according to the least such $r$ gives at most $M$ pieces, each vanishing where $j$ is not pivotal for a regular
function, so the Cauchy--Schwarz inequality and the proof of \cref{thm:multi} give
$\nu_{j,s}(F^j_{l_s}>v)\le MBe^{-v/2}$, while the bound on $\nu_{j,s}(l_s\text{ is not dense})$, which uses only
\cref{lem:level}, is unchanged. Hence $\E_{\nu_j}G_j\le\frac{4a}{\ln2}\bigl(\ln(2kMB)+1\bigr)+\frac2\gamma(\log k+\log e)$.

The path lemma (\cref{lem:upath}) holds for arbitrary $h_l$. The constant becomes
$53.7+12.75\log k+(1+4a)\log M$ with $a=1.44$. For $H^+$ the proof of \cref{thm:plus}(d) applies with the same
changes. In its second case $\Inf_j(h_{l_s})>1/(2k)$, so some $h_{l_s,r}$ has $\Inf_j(h_{l_s,r})>1/(2kM)$; since the
influences of a regular function do not increase along its order and their squares sum to at most $1$, this gives
$r_{l_s,r}(j)<4k^2M^2$ and $G_j<a\log(4k^2M^2)+\log(4k)$ pointwise. The constant becomes
$92.7+23.05\log k+(1+8a)\log M$.
\end{proof}

\section{The strengthened inequality and composition}\label{sec:plus}

Recall $H^+(f)=\sum_{A\ne\emptyset}P(A)\log(1/P(A))$ from \eqref{eq:feiplus}, and put
\[
W(f)=\sum_{|A|\ge2}P(A),\qquad J(f)=\sum_{|A|\ge2}|A|\,P(A)=\sum_{i=1}^n\Var(D_ig).
\]
The second expression for $J(f)$ holds because $\sum_{A\ni i,\,A\ne\{i\}}P(A)=\Inf_i(f)-\widehat{D_ig}(\emptyset)^2
=\Var(D_ig)$; note that $\widehat{D_ig}(\emptyset)=\hatg(\{i\})$. Since $|A|\le2(|A|-1)$ when $|A|\ge2$, we have
$J(f)\le2\bigl(I(f)-\Var(f)\bigr)$, so it suffices to bound $H^+(f)$ by a multiple of $J(f)$. The next two lemmas
hold for every Boolean function.

\begin{lemma}\label{lem:single}
For every $i\in[n]$, $\;P(\{i\})\log\dfrac1{P(\{i\})}\le\dfrac2{\ln2}\,\Var(D_ig)$.
\end{lemma}

\begin{proof}
Let $b=|\hatg(\{i\})|=|\E D_ig|$. Then $b\le\E|D_ig|=\Inf_i(f)$, so $\Var(D_ig)=\Inf_i(f)-b^2\ge b-b^2$. Since
$\ln(1/b)\le1/b-1$, we get $P(\{i\})\log(1/P(\{i\}))=\frac2{\ln2}\,b^2\ln\frac1b\le\frac2{\ln2}(b-b^2)$.
\end{proof}

The constant in \cref{lem:single} cannot be improved: for $f=x_1\wedge\neg(x_2\wedge\cdots\wedge x_w)$ and $i=1$
the ratio of the two sides tends to $2/\ln 2$ as $w\to\infty$.

\begin{lemma}\label{lem:mass}
$W(f)\log\dfrac1{W(f)}\le8\,J(f)$.
\end{lemma}

\begin{proof}
Write $W=W(f)$, $J=J(f)$, $W_1=\sum_iP(\{i\})$ and $\alpha=\min(p,1-p)$. We use four facts.
(i) $J\ge\sum_i\Inf_i(f)\bigl(1-\Inf_i(f)\bigr)$, since $\Var(D_ig)=\Inf_i(f)-\hatg(\{i\})^2$ and
$|\hatg(\{i\})|\le\Inf_i(f)$.
(ii) $W\le J/2$ and $W\le\Var(f)\le4\alpha$.
(iii) Every edge in direction $i$ along which $f$ changes has an endpoint in the smaller of $f^{-1}(0)$ and
$f^{-1}(1)$, so $\Inf_i(f)\le2\alpha$. Hence $W_1\le\sum_i\Inf_i(f)^2\le2\alpha\,I(f)$ and
$J=I(f)-W_1\ge(1-2\alpha)I(f)$.
(iv) $I(f)\ge2\alpha\log(1/\alpha)$.
Recall also that $x\log(1/x)$ is increasing on $(0,1/e]$ and never exceeds $(\log e)/e<0.531$.

\emph{Case $\alpha\le1/4$.} By (iii) and (iv), $J\ge I(f)/2\ge\alpha\log(1/\alpha)$. If $4\alpha\le1/e$, then
(ii) gives $W\log(1/W)\le4\alpha\log\frac1{4\alpha}\le4J$. If $4\alpha>1/e$, then $J\ge\alpha\log(1/\alpha)\ge0.316$,
so $W\log(1/W)<0.531<1.7\,J$.

\emph{Case $\alpha>1/4$ and $J\ge\delta:=0.024$.} Put $w_0=0.012$. If $W\ge w_0$, then (ii) gives
$W\log(1/W)\le\frac J2\log\frac1{w_0}<3.2\,J$. If $W<w_0$, then $W\log(1/W)\le w_0\log(1/w_0)<3.2\,\delta\le3.2\,J$.

\emph{Case $\alpha>1/4$ and $J<\delta$.} Now $\Var(f)>3/4$, so $I(f)\ge\Var(f)>3/4$, while $I(f)=W_1+J<1+\delta$.
By (i), every coordinate has $\Inf_i(f)(1-\Inf_i(f))<\delta$, hence $\Inf_i(f)<2\delta$ or $\Inf_i(f)>1-2\delta$.
Coordinates of the first kind satisfy $\Inf_i(f)\le\Inf_i(f)(1-\Inf_i(f))/(1-2\delta)$, so their total influence is
less than $\delta/(1-2\delta)<3/4$. Therefore a coordinate $i^\ast$ of the second kind exists, and it is unique
because $2(1-2\delta)>1+\delta$. After negating $x_{i^\ast}$ if necessary, $b:=\hatg(\{i^\ast\})=\E D_{i^\ast}g$ satisfies
$b^2=\Inf_{i^\ast}(f)-\Var(D_{i^\ast}g)>1-3\delta$, so $b>1-3\delta$.

Write $g=A+\chi_{\{i^\ast\}}B$ with $A=\E_{i^\ast}g$ and $B=D_{i^\ast}g$. These are functions of the other
variables with values in $\{-1,0,1\}$, and exactly one of them vanishes at each point. We have $\hatg(T)=\widehat A(T)$
for $T\not\ni i^\ast$ and $\hatg(T\cup\{i^\ast\})=\widehat B(T)$. Let $D=1-B\in\{0,1,2\}$ and let $E=\{D\ne0\}$, of
density $\eta$. From $b=1-\E D$ and $\E D\ge\eta$ we get $\eta\le1-b<3\delta$. The sets with at least two elements
are the sets $T\not\ni i^\ast$ with $|T|\ge2$ and the sets $T\cup\{i^\ast\}$ with $T\ne\emptyset$; since
$\|A\|_2^2=\Pr[B=0]=\Pr[D=1]$,
\[
W\le\|A\|_2^2+\Var(B)\le\Pr[D=1]+\E[D^2]\le4\eta .
\]
On the other hand $J\ge\sum_{T\ne\emptyset}(|T|+1)\widehat B(T)^2\ge\sum_T|T|\widehat D(T)^2=\sum_j\E[(D_jD)^2]$.
If an edge in direction $j$ leaves $E$, then $|D_jD|\ge\frac12$ at both of its endpoints. Hence
$\sum_j\E[(D_jD)^2]$ is at least $\frac14$ of the total influence of the indicator of $E$, which is at least
$2\eta\log(1/\eta)$ by edge isoperimetry. Thus $J\ge\frac\eta2\log\frac1\eta$, and since $4\eta<12\delta<1/e$,
\[
W\log\frac1W\le4\eta\log\frac1{4\eta}\le4\eta\log\frac1\eta\le8J.\qedhere
\]
\end{proof}

\begin{proposition}\label{prop:reduce}
Let $f$ be a Boolean function and $Q\ge0$ a function on the subsets of $[n]$ with $\sum_{|A|\ge2}Q(A)\le1$ and
$Q(A)>0$ whenever $|A|\ge2$ and $P(A)>0$. Then
\[
H^+(f)\le\Bigl(\frac2{\ln2}+8\Bigr)J(f)+\sum_{|A|\ge2}P(A)\log\frac1{Q(A)} .
\]
\end{proposition}

\begin{proof}
Split $H^+(f)$ into the contribution of the singletons and that of the sets with at least two elements. By
\cref{lem:single}, the first is at most $\frac2{\ln2}\sum_i\Var(D_ig)=\frac2{\ln2}J(f)$. If $W=W(f)>0$, Gibbs'
inequality applied to the probability distribution $P(A)/W$ on $\{A:|A|\ge2\}$ gives
$\sum_{|A|\ge2}P(A)\log\frac1{P(A)}\le W\log\frac1W+\sum_{|A|\ge2}P(A)\log\frac1{Q(A)}$, and \cref{lem:mass} bounds
$W\log(1/W)$.
\end{proof}

\begin{theorem}\label{thm:plus}
\textup{(a)} If $f$ is computed by a read-$k$ DNF or CNF formula, then $H^+(f)\le\min(10.9+14k,\,44+8.2\log k)\,J(f)$.

\noindent\textup{(b)} If $f$ is regular, then $H^+(f)\le30\,J(f)$.

\noindent\textup{(c)} If $f=F(h_1,\dots,h_m)$ as in \cref{thm:multi}, then $H^+(f)\le(33+15.3\ln m)\,J(f)$.

\noindent\textup{(d)} If $f$ is a union or an intersection of regular functions as in \cref{thm:union}, then
$H^+(f)\le\min(11.9+33k,\,93+23.1\log k)\,J(f)$.
\end{theorem}

Since $J(f)\le2(I(f)-\Var(f))$, \cref{thm:plus} implies \cref{thm:plus-intro}.

\begin{proof}[Proof of \cref{thm:plus}\textup{(a)}]
By duality, as in the proof of \cref{thm:readk}, it suffices to treat DNF formulas. Take $Q=Q_1$ from the proof of
\cref{thm:readk-main}, a probability measure, and $\beta=3$. For $|A|\ge2$ and $i\in A$ the set $A\setminus\{i\}$
is nonempty, so the cover used in that proof gives
\[
\sum_{|A|\ge2}P(A)\log\frac1{Q_1(A)}\le W(f)\,\Lambda+2\beta k\,J(f)+2\sum_i\Inf_i(f)\,w_i\,\nu_i(E_i\setminus\{\emptyset\}).
\]

\emph{The vacuum term.} Recall $\Lambda\le2\mu$. If $p_0\ge1/4$, then $\mu\le k\,I(f)$ by \cref{lem:path}, and
$W(f)I(f)=W(f)W_1+W(f)J(f)\le\frac32J(f)$ by fact (ii) in the proof of \cref{lem:mass}, so $W(f)\Lambda\le3k\,J(f)$.
If $p_0<1/4$, then $W(f)\le\Var(f)\le4p_0$ and $J(f)\ge I(f)/2$ by fact (iii), so
$W(f)\Lambda\le8p_0\mu\le2k\,I(f)\le4k\,J(f)$.

\emph{Sparse elements.} We show $w_i\,\nu_i(E_i\setminus\{\emptyset\})\le\max(\log k+2,\,2kC_\beta)\,\nu_i(T\ne\emptyset)$,
where $\nu_i(T\ne\emptyset)=\Var(D_ig)/\Inf_i(f)$. If $w_i<\log k+2$ this is clear. Otherwise every term through $i$
has at least $\log k+2$ variables. In the notation of the proof of \cref{lem:witness},
$\widehat{d_s}(\emptyset)=2^{-|W_s|}\widehat{\psi_s}(\emptyset)$ and $\|d_s\|_2^2=2^{-|W_s|}\|\psi_s\|_2^2$, so
$|\widehat{d_s}(\emptyset)|\le2^{-|W_s|/2}\|d_s\|_2$, and the Cauchy--Schwarz inequality gives
\[
\nu_i(\emptyset)=\frac{\widehat{D_ig}(\emptyset)^2}{\Inf_i(f)}\le\sum_s2^{-|W_s|}\le k\,2^{-\log k-1}=\frac12 .
\]
Hence $\nu_i(T\ne\emptyset)\ge\frac12$, while $w_i\nu_i(E_i)\le kC_\beta$ as in the proof of \cref{thm:readk-main}.
Summing with weights $\Inf_i(f)$, the sparse term is at most $2\max(\log k+2,\,2kC_\beta)\,J(f)$.

With $\beta=3$ and $C_3<0.949$, the three terms add up to at most $\bigl(4k+6k+2\max(\log k+2,1.898k)\bigr)J(f)\le14k\,J(f)$,
and \cref{prop:reduce} gives $H^+(f)\le(2/\ln2+8+14k)\,J(f)$.

For the logarithmic bound take $Q=Q_1$ from the proof of \cref{thm:readk-log}. The bound on $\log(1/Q_1(A))$ obtained
there holds for every $A\ne\emptyset$ and has no term independent of $|A|$ other than $\Lambda_\lambda$, so
\[
\sum_{|A|\ge2}P(A)\log\frac1{Q_1(A)}\le W(f)\Lambda_\lambda+K_\beta\,J(f)+\kappa_\beta\sum_i\Inf_i(f)\,w_i\,\nu_i(E_i\setminus\{\emptyset\}),
\]
with $K_\beta=\log k+\beta+1+2\log3+2\log e+\log(\beta k+1)+\log(\beta k)$. As in the vacuum step above, with
$\Lambda_\lambda\le(\log e)\mu/k$ in place of $\Lambda\le2\mu$: if $p_0\ge\frac14$ then $W(f)\Lambda_\lambda\le(\log e)
W(f)I(f)\le\frac32(\log e)J(f)$, and if $p_0<\frac14$ then $W(f)\Lambda_\lambda\le4p_0(\log e)\mu/k\le(\log e)I(f)\le
2(\log e)J(f)$. The sparse term is at most $\max(\log k+2,\,2M_\beta(k))\,J(f)$ by the argument above, with
$w_i\nu_i(E_i)\le M_\beta(k)$ from the proof of \cref{thm:readk-log} in place of $kC_\beta$. Hence
\[
H^+(f)\le\Bigl(\frac2{\ln2}+8+2\log e+K_\beta+\kappa_\beta\max\bigl(\log k+2,\,2M_\beta(k)\bigr)\Bigr)J(f),
\]
and $\beta=8$ with \eqref{eq:Mbound} gives at most $35.00+8.48<44$ for $k=1$ and at most
$37.35+8.15\log k$ for $k\ge2$.
\end{proof}

\begin{proof}[Proof of \cref{thm:plus}\textup{(b)}]
As in \cref{sec:regular}, assume that $f$ is monotone with coordinates ordered by decreasing desirability. Then
$D_ig\ge0$, so $\hatg(\{i\})=\Inf_i(f)$ and $\nu_i(\emptyset)=\Inf_i(f)$. Influences decrease along the order: if
$i<j$ and $z$ denotes the remaining coordinates, then
$\Inf_i(f)-\Inf_j(f)=\E_z\bigl[f(x_i{=}1,x_j{=}0,z)-f(x_i{=}0,x_j{=}1,z)\bigr]\ge0$. Since
$\sum_i\Inf_i(f)^2=\sum_iP(\{i\})\le1$, a coordinate with $\Inf_i(f)>1/2$ has $i\le3$.

We claim that
\[
\sum_{|A|\ge2}P(A)\sum_{i\in A}\ln^+\frac i{|A|}\le2L_1\,J(f).
\]
The left side equals $\sum_i\Inf_i(f)\,\E_{\nu_i}\bigl[\ln^+\frac i{|T|+1}\,\mathbf 1[T\ne\emptyset]\bigr]$, and
$\ln^+\frac i{|T|+1}\le\ln\frac i{|T\cap U_i|+1}$. If $\Inf_i(f)\le\frac12$, then $\nu_i(T\ne\emptyset)\ge\frac12$
and \cref{lem:ranktail} bounds the expectation by $L_1\le2L_1\,\nu_i(T\ne\emptyset)$. If $\Inf_i(f)>\frac12$, then
$i\le3$ and the expectation is at most $(\ln3)\,\nu_i(T\ne\emptyset)$. Summing with weights $\Inf_i(f)$ proves the
claim.

Now let $Q(A)=2^{-(|A|-1)}\pi_{|A|}(A)$ for $|A|\ge2$, with $\pi_k$ as in the proof of \cref{lem:rankref}. Then
$\sum_{|A|\ge2}Q(A)=1$ and, by the bound on $Z_k$ in that proof,
$\log\frac1{Q(A)}\le|A|\bigl(1+\log\frac{ea}{a-1}\bigr)+\frac a{\ln2}\sum_{i\in A}\ln^+\frac i{|A|}$. By
\cref{prop:reduce} and the claim,
\[
H^+(f)\le\Bigl(\frac2{\ln2}+9+\log\frac{ea}{a-1}+\frac{2aL_1}{\ln2}\Bigr)J(f),
\]
and at $a=1.1087$ the bracket is less than $30$.
\end{proof}

\begin{proof}[Proof of \cref{thm:plus}\textup{(c)}]
Use the notation of the proof of \cref{thm:multi}. We claim that for every $i$
\[
\E_{\nu_i}\bigl[G_i\,\mathbf 1[T\ne\emptyset]\bigr]\le(8\ln m+9.86)\,\nu_i(T\ne\emptyset).
\]
If $\nu_i(T\ne\emptyset)\ge\frac12$, this follows from $\E_{\nu_i}G_i<4\ln m+4.93$. Otherwise
$\nu_i(\emptyset)>\frac12$, and since $\nu_i(\emptyset)=\hatg(\{i\})^2/\Inf_i(f)\le\Inf_i(f)$ we get $\Inf_i(f)>\frac12$.
As $\Inf_i(f)\le\sum_l\Inf_i(h_l)$, some $h_l$ has $\Inf_i(h_l)>1/(2m)$. The influences of the regular function $h_l$
do not increase along its order (see the proof of part (b)) and their squares sum to at most $1$, so
$r_l(i)<4m^2$, and then $G_i\le F_l\le\ln r_l(i)<2\ln(2m)$ pointwise. This proves the claim, and summing it with
weights $\Inf_i(f)$ gives
\[
\sum_{|A|\ge2}P(A)\sum_{i\in A}\min_l\ln^+\frac{r_l(i)}{|A|}\le(8\ln m+9.86)\,J(f).
\]
Now take $Q(A)=2^{-(|A|-1)}\pi_{|A|}(A)$ with the weights $\theta_k$ of the proof of \cref{thm:multi}; as in part (b),
\cref{prop:reduce} gives
\[
H^+(f)\le\Bigl(\frac2{\ln2}+9+\log\frac{ema}{a-1}+\frac{a(8\ln m+9.86)}{\ln2}\Bigr)J(f),
\]
and $a=1.2$ gives the stated bound.
\end{proof}

\begin{proof}[Proof of \cref{thm:plus}\textup{(d)}]
By complementation it suffices to treat unions. Take $Q=Q_1$ and the constants of the proof of \cref{thm:union}.
The bound on $\log(1/Q_1(A))$ obtained there holds for every $A\ne\emptyset$, so
\[
\sum_{|A|\ge2}P(A)\log\frac1{Q_1(A)}\le W(f)\Lambda+(\beta+c_1)k\,J(f)+J(f)
+\sum_j\Inf_j(f)\,\E_{\nu_j}\bigl[G_j\,\mathbf 1[T\ne\emptyset]\bigr].
\]
The vacuum term is treated as in part (a), with \cref{lem:upath} in place of \cref{lem:path}:
$W(f)\Lambda\le4k\,J(f)$. We claim that
$\E_{\nu_j}[G_j\mathbf 1[T\ne\emptyset]]\le2k\bigl(\frac{aL_1}{\ln2}+c_2K\bigr)\nu_j(T\ne\emptyset)$. If
$\nu_j(T\ne\emptyset)\ge\frac12$, this follows from the bound on $\E_{\nu_j}G_j$ in the proof of \cref{thm:union}.
Otherwise $\Inf_j(f)\ge\nu_j(\emptyset)>\frac12$ as in part (c); since $\Inf_j(f)\le\sum_s\Inf_j(h_{l_s})$ with at most
$k$ summands, some $s$ has $\Inf_j(h_{l_s})>1/(2k)$. Then $r_{l_s}(j)<4k^2$ as in part (c), and
$\Pr[h_{l_s}=1]\ge\Inf_j(h_{l_s})/2>1/(4k)$, so pointwise
$G_j\le G_j(l_s)<\frac a{\ln2}\ln(4k^2)+c_2\log(4k)\le2k\bigl(\frac{aL_1}{\ln2}+c_2K\bigr)$. Summing the claim with
weights $\Inf_j(f)$ and using \cref{prop:reduce},
\[
H^+(f)\le\Bigl(\frac2{\ln2}+9+\bigl(4+\beta+c_1+2\,\tfrac{aL_1}{\ln2}+2c_2K\bigr)k\Bigr)J(f),
\]
and the coefficient of $k$ is $4+5+4.273+17.248+2.457<33$.

For the logarithmic bound take $Q=Q_1$ and the constants of the proof of \cref{thm:union-log}. The bound on
$\log(1/Q_1(A))$ obtained there holds for every $A\ne\emptyset$ and has no term independent of $|A|$ other than
$\Lambda_\lambda$, so
\[
\sum_{|A|\ge2}P(A)\log\frac1{Q_1(A)}\le W(f)\Lambda_\lambda+K_u\,J(f)+\sum_j\Inf_j(f)\,\E_{\nu_j}\bigl[G_j\,\mathbf
1[T\ne\emptyset]\bigr],
\]
with $K_u=(1+a)\log k+\beta+1+2\log3+\log\frac{ea}{a-1}+\log(2e)$. As in part (a), with $\Lambda_\lambda\le(\log e)
\sum_l\pi_l/k$ and \cref{lem:upath}, $W(f)\Lambda_\lambda\le2(\log e)J(f)$. If $\nu_j(T\ne\emptyset)\ge\frac12$, then
$\E_{\nu_j}[G_j\mathbf1[T\ne\emptyset]]\le2E_k\,\nu_j(T\ne\emptyset)$, where $E_k=\frac{4a}{\ln2}(\ln(2kB)+1)+\frac2\gamma
(\log k+\log e)$ is the bound on $\E_{\nu_j}G_j$ from that proof; otherwise, as above, $G_j<a\log(4k^2)+\log(4k)$
pointwise. By \cref{prop:reduce},
\[
H^+(f)\le\Bigl(\frac2{\ln2}+8+2\log e+K_u+\max\bigl(2E_k,\,a\log(4k^2)+\log(4k)\bigr)\Bigr)J(f)
\le(93+23.1\log k)\,J(f).
\]
\end{proof}

We recall the biased form of \eqref{eq:feiplus}~\cite{OT13}. Let $\mu$ be a product probability measure on
$\bits^m$ with $q_j=\mu(y_j=1)\in(0,1)$, put $\nu_j=2q_j-1$, $\sigma_j^2=1-\nu_j^2=4q_j(1-q_j)$ and
$\phi_j(y)=(2y_j-1-\nu_j)/\sigma_j$, so that the products $\phi_A=\prod_{j\in A}\phi_j$ form an orthonormal basis of
$L^2(\mu)$. For $G:\bits^m\to\bits$ expand $1-2G=\sum_A\widetilde g(A)\phi_A$ and put $P_\mu(A)=\widetilde g(A)^2$,
$\sigma_A^2=\prod_{j\in A}\sigma_j^2$, $I_\mu(G)=\sum_A|A|P_\mu(A)$ and $\Var_\mu(G)=\sum_{A\ne\emptyset}P_\mu(A)$.
Then $G$ satisfies the $\mu$-biased form of \eqref{eq:feiplus} with constant $C$ if
\begin{equation}\label{eq:biased}
\sum_{A\ne\emptyset}P_\mu(A)\log\frac{\sigma_A^2}{P_\mu(A)}\le C\bigl(I_\mu(G)-\Var_\mu(G)\bigr).
\end{equation}
For the uniform measure this is \eqref{eq:feiplus}. Negating an input or the output of $G$ leaves every
$P_\mu(A)$ unchanged, up to replacing $q_j$ by $1-q_j$.

\begin{lemma}\label{lem:andor}
For every $m$ and every product measure $\mu$ on $\bits^m$, the functions $\mathrm{OR}_m$ and $\mathrm{AND}_m$
satisfy \eqref{eq:biased} with $C=2+4\log e=7.7707\ldots$, and $\mathrm{XOR}_m$ satisfies \eqref{eq:biased} with
$C=2\log e$. The constant $2+4\log e$ cannot be lowered: for $\mathrm{OR}_2$ with $q_1=q_2\to0$ the ratio of the two
sides of \eqref{eq:biased} tends to it.
\end{lemma}

\begin{proof}
First let $a_1,\dots,a_m\in(0,1]$, $Z=\prod_ja_j$, and
\[
\Phi(a)=\sum_j(a_j-Z)\ln\frac1{a_j},\qquad R(a)=\sum_j(1-a_j)-1+Z .
\]
We claim $\Phi\le2R$. For $m=1$ both sides vanish. Appending a number $b\in(0,1]$ replaces $Z$ by $bZ$ and gives
\[
\Phi(a,b)-\Phi(a)=(1-b)\,Z\ln\frac1Z+b\ln\frac1b\,(1-Z),\qquad R(a,b)-R(a)=(1-b)(1-Z),
\]
and since $t\ln(1/t)\le1-t$ for $0<t\le1$, the first difference is at most twice the second. The claim follows by
induction on $m$.

For $\mathrm{OR}_m$ it suffices, after negating the output, to take $G(y)=\prod_j(1-y_j)$, the indicator that all
inputs vanish. Put $a_j=1-q_j$. Since $1-y_j=a_j-\frac{\sigma_j}2\phi_j$,
\[
1-2G=1-2\prod_j\Bigl(a_j-\frac{\sigma_j}2\phi_j\Bigr),\qquad
P_\mu(A)=4\prod_{j\in A}\frac{\sigma_j^2}4\prod_{j\notin A}a_j^2\quad(A\ne\emptyset),
\]
and $P_\mu(\emptyset)=(1-2Z)^2$. Hence $\log(\sigma_A^2/P_\mu(A))=2|A|-2+2\sum_{j\notin A}\log(1/a_j)$ for
$A\ne\emptyset$. Using $\frac{\sigma_j^2}4+a_j^2=a_j$ one finds $\sum_{A\ne\emptyset,\,A\not\ni j}P_\mu(A)=4a_jZ-4Z^2$,
$I_\mu(G)=\sum_j\sigma_j^2\prod_{i\ne j}a_i=4Z\sum_j(1-a_j)$ and $\Var_\mu(G)=4Z(1-Z)$, so that
$I_\mu(G)-\Var_\mu(G)=4ZR$ and the left side of \eqref{eq:biased} equals
\[
2\bigl(I_\mu(G)-\Var_\mu(G)\bigr)+8Z\sum_j(a_j-Z)\log\frac1{a_j}
=8ZR+\frac{8Z\,\Phi}{\ln2}\le\Bigl(2+\frac4{\ln2}\Bigr)\bigl(I_\mu(G)-\Var_\mu(G)\bigr).
\]
$\mathrm{AND}_m$ is obtained from $\mathrm{OR}_m$ by negating the inputs and the output. For $m=2$ and $q_1=q_2=q$
we have $\Phi=2q(1-q)\ln\frac1{1-q}\sim2q^2$ and $R=q^2$, so the ratio of the two sides of \eqref{eq:biased} is
$2+2\Phi/(R\ln2)\to2+4\log e$ as $q\to0$.

For $\mathrm{XOR}_m$, $1-2G=\pm\prod_j(\nu_j+\sigma_j\phi_j)$, so $P_\mu(A)=\sigma_A^2\prod_{j\notin A}\nu_j^2$ and
$\log(\sigma_A^2/P_\mu(A))=\sum_{j\notin A}\log(1/\nu_j^2)$. By continuity we may assume $\nu_j\ne0$. With $a_j=\nu_j^2$,
$\sum_{A\not\ni j}P_\mu(A)=a_j$ and $P_\mu(\emptyset)=Z$, so the left side of \eqref{eq:biased} is $\Phi/\ln2$, while
$I_\mu(G)-\Var_\mu(G)=\sum_j\sigma_j^2-(1-Z)=R$.
\end{proof}

\subsection*{Threshold gates with inputs of equal bias}
For $1\le t\le r$ let $T_{t,r}(y)=\mathbf 1[y_1+\cdots+y_r\ge t]$ on $\bits^r$, and let $\mu$ be the product measure
with $\mu(y_j=1)=q$ for every $j$, where $0<q<1$. Write $N(y)=y_1+\cdots+y_r$, $N_{-j}=N-y_j$, $\sigma^2=4q(1-q)$
and $\Lambda=r\sigma^2$, and let $P=P_\mu$, $I=I_\mu(T_{t,r})$, $V=\Var_\mu(T_{t,r})$ and
$W_k=\sum_{|A|=k}P(A)$. Every coordinate is pivotal with the same probability $\rho=\mu(N_{-j}=t-1)$, and
$I=\Lambda\rho$.

\begin{theorem}\label{thm:threshold}
In this setting
\[
\sum_{A\ne\emptyset}P(A)\ln\frac{\sigma_A^2}{P(A)}\le 11.49\,I .
\]
Consequently $T_{t,r}$ satisfies \eqref{eq:biased} under $\mu$ with $C=66.4$, for every $r$, every $1\le t\le r$ and
every $q$.
\end{theorem}

The proof rests on two lemmas. The first is the analogue of \cref{lem:ranktail}.

\begin{lemma}\label{lem:thr-rank}
$\sum_{k\ge1}kW_k\ln^+(\Lambda/k)\le L'I$, where $L'=2\ln\bigl(3\sqrt{6/\pi}\bigr)+2<4.85$.
\end{lemma}

\begin{proof}
Let $d(y')=T_{t,r}(y',1)-T_{t,r}(y',0)=\mathbf 1[N(y')=t-1]$ for $y'\in\bits^{r-1}$, and let $\widetilde d(B)$ be
its coefficients in the basis $\phi_B$ of $L^2(\mu)$ on $\bits^{r-1}$. Since $\E_{y_r}[h\,\phi_r]=\frac\sigma2
(h(1)-h(0))$ for every function $h$ of $y_r$, we have $P(B\cup\{r\})=\sigma^2\widetilde d(B)^2$, and by symmetry
$\sum_{|A|=k,\,A\ni r}P(A)=\frac krW_k$. Hence $kW_k=\Lambda\sum_{|B|=k-1}\widetilde d(B)^2$, and
$\nu(B)=\widetilde d(B)^2/\rho$ is a probability distribution because $\E d^2=\rho$.

Let $y'$ have law $\mu$, fix $0<\eta\le1$, and let $z$ be obtained from $y'$ by replacing each coordinate,
independently with probability $\eta$, by a fresh sample from $\mu$. Then $\E[\phi_B(z)\mid y']=(1-\eta)^{|B|}
\phi_B(y')$, so $\E[d(y')d(z)]=\rho\,\E_\nu(1-\eta)^{|T|}$, while $\E[d(y')d(z)]\le\rho\max_{y'}\Pr[N(z)=t-1\mid y']$.
Write each replaced coordinate as $z_j=u_j\wedge b_j$ if $q\le\frac12$ and as $z_j=u_j\vee b_j$ if $q>\frac12$, with
$u_j$ a uniform bit and $b_j$ an independent bit equal to $1$ with probability $2q$, respectively $2q-1$. Given the
set of replaced coordinates and the bits $b_j$, the replaced coordinates with $b_j=1$ (respectively $b_j=0$) are
independent uniform bits, say $F$ of them, and every other coordinate of $z$ is fixed; so by \cref{lem:sperner}
the conditional probability that $N(z)=t-1$ is at most $\frac2{\sqrt\pi}(F+1)^{-1/2}$. The number $F$ has the law
$\Bin(r-1,2\eta\tilde q)$ with $\tilde q=\min(q,1-q)\ge q(1-q)$, so $\E[1/(F+1)]\le1/(2r\eta\tilde q)\le2/(\eta\Lambda)$
and, by Jensen's inequality, $\E[(F+1)^{-1/2}]\le(2/(\eta\Lambda))^{1/2}$. For $K\ge1$ take $\eta=1/(3K)$; as in the
proof of \cref{lem:ranktail},
\[
\nu(|T|\le K)\le\Bigl(1-\frac1{3K}\Bigr)^{-K}\E_\nu(1-\eta)^{|T|}\le3\sqrt{6/\pi}\,\Bigl(\frac K\Lambda\Bigr)^{1/2}.
\]
Finally $\sum_kkW_k\ln^+(\Lambda/k)=I\,\E_\nu\ln^+\frac\Lambda{|T|+1}$, and the tail bound with $K=\Lambda e^{-u}$ gives
\[
\E_\nu\ln^+\frac\Lambda{|T|+1}\le\int_0^\infty\min\bigl(1,3\sqrt{6/\pi}\,e^{-u/2}\bigr)\,du=L'.
\]
\end{proof}

The second lemma pays for the normalisation $V\ln(1/V)$. When the inputs are nearly always $0$ and $t\ge2$ this
term is much larger than $I$, and it is repaid by the factors $\sigma_A^2$; the lemma measures this through the ratio
$\rho/V$.

\begin{lemma}\label{lem:thr-credit}
$V\ln(1/V)\le I\ln^+(4\rho/V)+C_2I$, where $C_2=1+\ln2+8/e<4.64$.
\end{lemma}

\begin{proof}
Let $\alpha=\min(\mu(T_{t,r}=1),\mu(T_{t,r}=0))$. If $\alpha\ge\frac14$ then $V\ge\frac34$ and
$V\ln(1/V)\le V\ln\frac43<0.29\,I$. Let $\alpha<\frac14$. Negating all inputs and the output maps $T_{t,r}$ to
$T_{r-t+1,r}$ and $q$ to $1-q$, and changes none of $V$, $I$, $\rho$; so we may assume $\alpha=\mu(A)$ with
$A=\{N\ge t\}$. Then $t>rq$, since otherwise $t\le\lfloor rq\rfloor$, which is at most the median of $\Bin(r,q)$
\cite{KB80}, and $\mu(A)\ge\frac12$. As $V=4\alpha(1-\alpha)\le4\alpha$ and $4(1-\alpha)>1$, we have
$V\ln(1/V)\le4\alpha\ln(1/\alpha)$.

Let $\mu_A=\mu(\cdot\mid A)$ and let $\mu_A^{(j)}$ be the law of $y_j$ under $\mu_A$. By the chain rule for relative
entropy with respect to a product measure,
\[
\ln\frac1\alpha=D(\mu_A\,\|\,\mu)=\mathrm{TC}(\mu_A)+\sum_jD\bigl(\mu_A^{(j)}\,\big\|\,\mu_j\bigr),\qquad
\mathrm{TC}(\mu_A)=D\Bigl(\mu_A\,\Big\|\,\textstyle\bigotimes_j\mu_A^{(j)}\Bigr).
\]
For a coordinate $j$ put $a=\mu(A\mid y_j=1)$ and $b=\mu(A\mid y_j=0)$, so that $\alpha=qa+(1-q)b$ and $a-b=\rho$.
Since $b\le\alpha$, $\alpha D(\mu_A^{(j)}\|\mu_j)=qa\ln\frac a\alpha+(1-q)b\ln\frac b\alpha\le qa\ln\frac a\alpha$.
With $u=(1-q)\rho/\alpha$ we have $a/\alpha=1+u$ and $(1+u)\ln(1+u)\le u(1+\ln(1+u))$, hence
\[
\alpha D\bigl(\mu_A^{(j)}\,\big\|\,\mu_j\bigr)\le q(1-q)\rho\,(1+\ln(1+u))\le\frac{\sigma^2\rho}4\Bigl(1+\ln2+\ln^+\frac{4\rho}V\Bigr),
\]
using $\ln(1+u)\le\ln2+\ln^+u$ and $\alpha\ge V/4$. Summing over $j$ gives $4\alpha\sum_jD(\mu_A^{(j)}\|\mu_j)\le
I(1+\ln2+\ln^+(4\rho/V))$.

It remains to show $\alpha\,\mathrm{TC}(\mu_A)\le\frac2eI$. If $t=r$ then $\mu_A$ is a point mass and
$\mathrm{TC}(\mu_A)=0$. Let $rq<t\le r-1$ and let $R$ be the product measure with $R(y_j=1)=t/r$; its density is
$dR/d\mu=M^{-1}e^{\lambda N}$ with $\lambda=\ln\frac{t(1-q)}{(r-t)q}>0$ and $M=\E_\mu e^{\lambda N}$. For every
product measure $R$ we have $D(\mu_A\|R)=\mathrm{TC}(\mu_A)+\sum_jD(\mu_A^{(j)}\|R_j)\ge\mathrm{TC}(\mu_A)$, and here
\[
D(\mu_A\,\|\,R)=\ln\frac M\alpha-\lambda\,\E_{\mu_A}N\le\ln\frac{Me^{-\lambda t}}\alpha,\qquad
\alpha=M\,\E_R\bigl[e^{-\lambda N}\mathbf 1_A\bigr]\ge Me^{-\lambda t}R(N=t).
\]
Hence $\mathrm{TC}(\mu_A)\le\ln(1/b^\ast)$, where $b^\ast=R(N=t)=\binom rt(t/r)^t(1-t/r)^{r-t}\ge(8v)^{-1/2}$ with
$v=t(r-t)/r\ge\frac12$~\cite[Ch.~10, Lemma~7]{MS77}. In the same way
$\rho=Me^{-\lambda(t-1)}R(N_{-j}=t-1)/(1-q+qe^\lambda)$ and $\alpha\le Me^{-\lambda t}$, and
$R(N_{-j}=t-1)=\binom{r-1}{t-1}(t/r)^{t-1}(1-t/r)^{r-t}=b^\ast$; since
$\sigma^2e^\lambda/(1-q+qe^\lambda)=4\frac tr(1-\frac tr)(1-q+qe^\lambda)\ge4\frac tr(1-\frac tr)$, this gives
$I/\alpha\ge4vb^\ast\ge(2v)^{1/2}$. Therefore
\[
\alpha\,\mathrm{TC}(\mu_A)\le\frac12\ln(8v)\,\frac I{(2v)^{1/2}}=\frac{\ln(8v)}{\sqrt{8v}}\,I\le\frac2e\,I,
\]
and altogether $V\ln(1/V)\le4\alpha\ln(1/\alpha)\le I\ln^+(4\rho/V)+(1+\ln2+8/e)I$.
\end{proof}

\begin{proof}[Proof of \cref{thm:threshold}]
By symmetry $P(A)=W_k/\binom rk$ when $|A|=k$, so the left side equals
\[
\sum_{k\ge1}W_k\ln\frac1{W_k}+\sum_{k\ge1}W_k\ln\Bigl(\binom rk\sigma^{2k}\Bigr).
\]
The first sum is $V\ln(1/V)+V\,H(K)$, where $K$ has law $W_k/V$ on $\{1,\dots,r\}$ and mean $m=I/V$. Gibbs'
inequality against the geometric law with mean $m$ gives $H(K)\le1+\ln m$, and $V(1+\ln(I/V))\le I$ because
$x(1+\ln(1/x))\le1$ for $0<x\le1$. For the second sum, $\binom rk\le(er/k)^k$ gives
$\ln(\binom rk\sigma^{2k})\le k+k\ln(\Lambda/k)$. If $4m\le\Lambda$, then $\ln^+(4\rho/V)=\ln^+(4m/\Lambda)=0$ and
$\sum_kkW_k\ln(\Lambda/k)\le L'I$ by \cref{lem:thr-rank}. If $4m>\Lambda$, write
$\ln\frac\Lambda k=-\ln\frac{4m}\Lambda+\ln\frac{4m}k$ and use $\ln x\le x/e$:
$\sum_kkW_k\ln(\Lambda/k)\le-I\ln(4\rho/V)+\frac{4m}e V$, and $4mV/e=4I/e<L'I$. In both cases the left side is at most
\[
V\ln\frac1V+2I-I\ln^+\frac{4\rho}V+L'I\le(2+L'+C_2)\,I<11.49\,I
\]
by \cref{lem:thr-credit}.

For the second statement, \cref{lem:andor} covers $t\in\{1,r\}$. Let $2\le t\le r-1$. The number
$\rho=\Pr[\Bin(r-1,q)=t-1]$ is the value of a binomial distribution at an interior point; it is $2q(1-q)\le\frac12$
when $r=3$, and for $r\ge4$ it is at most $\max_{1\le i\le m-1}\binom mi(i/m)^i(1-i/m)^{m-i}\le(m/(2\pi(m-1)))^{1/2}
<\frac12$ with $m=r-1$~\cite{MS77}. Since $d\ge0$ we have $P(\{j\})=\sigma^2(\E d)^2=\sigma^2\rho^2$ and
$\sum_{A\ni j}P(A)=\sigma^2\rho$, so $\sum_{|A|\ge2}|A|P(A)=I-\Lambda\rho^2\ge I/2$ and
$I-V=\sum_{|A|\ge2}(|A|-1)P(A)\ge I/4$. In bits the left side of \eqref{eq:biased} is therefore at most
$\frac{4\cdot11.49}{\ln2}(I-V)<66.4\,(I-V)$.
\end{proof}

\begin{remark}\label{rem:threshold}
The same reduction shows where general Boolean functions behave differently. For the disjunction of $s$ disjoint
conjunctions of two inputs, each input equal to $1$ with probability $s^{-2}$, the total correlation of $\mu_A$ grows
like $\ln s$ (it records which conjunction is satisfied) while $I_\mu/\alpha$ stays bounded, and the inequality of
\cref{lem:thr-credit}, with $\sum_j\Inf_j\ln^+(4\rho_j/V)$ in place of $I\ln^+(4\rho/V)$, fails for these functions.
\end{remark}

\subsection*{Threshold gates under arbitrary product measures}
Under an arbitrary product measure the same argument proves the inequality with $I_\mu$ in place of
$I_\mu-\Var_\mu$ on the right. For read-once formulas this weaker form costs a factor of the depth.

\begin{proposition}\label{prop:thr-general}
Let $1\le t\le r$ and let $\mu$ be a product measure on $\bits^r$ with $\mu(y_j=1)=q_j\in(0,1)$. Then
\[
\sum_{A\ne\emptyset}P_\mu(A)\ln\frac{\sigma_A^2}{P_\mu(A)}\le26.3\,I_\mu(T_{t,r}).
\]
\end{proposition}

\begin{proof}
Write $N$, $N_{-j}$ and $W_k$ as before, $\sigma_j^2=4q_j(1-q_j)$, $\Lambda=\sum_j\sigma_j^2$,
$\rho_j=\mu(N_{-j}=t-1)$, $\Inf_j=\sigma_j^2\rho_j$, $I=\sum_j\Inf_j$, $V=\Var_\mu(T_{t,r})$, $m=I/V$ and
$\pi_j=\Inf_j/V$, so that $\sum_j\pi_j=m$.

For $k\ge1$ let $\theta_k(j)=\max\bigl(\sigma_j^2\min(1,k/\Lambda),\,4\pi_j\bigr)$ and let $Z_k$ be the $k$-th elementary
symmetric polynomial of $\theta_k(1),\dots,\theta_k(r)$. Gibbs' inequality on the sets of size $k$, against the weights
$\prod_{j\in A}\theta_k(j)$, gives
\[
\sum_{|A|=k}P(A)\ln\frac{\sigma_A^2}{P(A)}\le W_k\ln\frac1{W_k}+W_k\ln Z_k+\sum_{|A|=k}P(A)\sum_{j\in A}\ln\frac{\sigma_j^2}{\theta_k(j)}.
\]
Since $\sum_j\theta_k(j)\le\min(\Lambda,k)+4m\le k+4m$, we have $Z_k\le(e(k+4m)/k)^k$ and $\ln Z_k\le k+4m$. Moreover
$\ln(\sigma_j^2/\theta_k(j))\le\min\bigl(\ln^+(\Lambda/k),\ln(V/(4\rho_j))\bigr)\le\ln^+(\Lambda/k)-\ln^+(4\rho_j/V)$.
Summing over $k$, and bounding $\sum_kW_k\ln(1/W_k)$ by $V\ln(1/V)+I$ as in the proof of \cref{thm:threshold},
\[
\sum_{A\ne\emptyset}P(A)\ln\frac{\sigma_A^2}{P(A)}\le V\ln\frac1V+6I+\sum_jI_j'-\sum_j\Inf_j\ln^+\frac{4\rho_j}V,
\qquad I_j'=\sum_{A\ni j}P(A)\ln^+\frac\Lambda{|A|}.
\]

For the rank terms, $I_j'=\Inf_j\,\E_{\nu_j}\ln^+\frac\Lambda{|T|+1}$ with $\nu_j(T)=P(T\cup\{j\})/\Inf_j$. If
$\Lambda\le2$ the expectation is at most $\ln2$. If $\Lambda>2$, the proof of \cref{lem:thr-rank} applies to the
coordinate $j$, except that a replaced coordinate $i$ is now free with probability $2\eta\tilde q_i\ge\eta\sigma_i^2/2$,
so that $F$ is a sum of independent indicators with mean at least $\eta(\Lambda-\sigma_j^2)/2\ge\eta\Lambda/4$; for
such a sum $\E[1/(F+1)]\le1/\E F$, hence $\nu_j(|T|\le K)\le3\sqrt{12/\pi}\,(K/\Lambda)^{1/2}$ and
$\E_{\nu_j}\ln^+\frac\Lambda{|T|+1}\le2\ln\bigl(3\sqrt{12/\pi}\bigr)+2<5.54$. So $\sum_jI_j'\le5.54\,I$.

It remains to show $V\ln(1/V)\le\sum_j\Inf_j\ln^+(4\rho_j/V)+14.7\,I$, which gives the constant $6+5.54+14.7<26.3$. We
follow the proof of \cref{lem:thr-credit}. If $\alpha\ge\frac14$ the bound is immediate. Otherwise, after negating all
inputs if necessary, $\alpha=\mu(A)$ with $A=\{N\ge t\}$, and $t>\E_\mu N$: if $t\le\E_\mu N=r\bar q$, then
$t\le\lfloor r\bar q\rfloor$, and Hoeffding's comparison with the binomial law~\cite[Theorem~4]{Hoe56} together
with~\cite{KB80} gives $\mu(A)\ge\Pr[\Bin(r,\bar q)\ge t]\ge\frac12$. The bound
$4\alpha D(\mu_A^{(j)}\|\mu_j)\le\Inf_j(1+\ln2+\ln^+(4\rho_j/V))$ holds for each coordinate as before. For the total
correlation, if $t=r$ it vanishes; otherwise let $R$ be the product measure with $dR/d\mu\propto e^{\lambda N}$, with
$\lambda>0$ chosen so that $\E_RN=t$. As in the proof of \cref{lem:thr-credit}, $\mathrm{TC}(\mu_A)\le\ln(1/R(N=t))$
and, for each $j$, $\sigma_j^2\rho_j/\alpha\ge4R(y_j{=}1)R(y_j{=}0)R(N_{-j}=t-1)$. Applying the identity
$\E_R[(y_j-R(y_j{=}1))h(N)]=R(y_j{=}1)R(y_j{=}0)\,\E_R[h(N_{-j}+1)-h(N_{-j})]$ to $h=\mathbf 1[\,\cdot\ge t]$ and summing over
$j$ gives $\sum_jR(y_j{=}1)R(y_j{=}0)R(N_{-j}=t-1)=\E_R(N-t)^+$, so $I/\alpha\ge4\E_R(N-t)^+=2\E_R|N-t|$. With
$v=\Var_R(N)$ we have $\E_R(N-t)^4\le v+3v^2$, and H\"older's inequality gives
$\E_R|N-t|\ge v^{3/2}(v+3v^2)^{-1/2}$; hence $I/\alpha\ge2v(1+3v)^{-1/2}$. By Darroch's theorem~\cite{Dar64}, $t$ is a mode
of $N$ under $R$. If $v\le\frac12$, Chebyshev's inequality gives $R(N=t)\ge1-v$ and $\mathrm{TC}(\mu_A)\le\ln\frac1{1-v}\le
2v\ln2$, so $\alpha\,\mathrm{TC}(\mu_A)\le(\ln2)(1+3v)^{1/2}I<1.1\,I$. If $v>\frac12$, let $K=\lceil2\sqrt v\rceil$; by
Chebyshev's inequality $R(|N-t|\le K)\ge\frac34$, so $R(N=t)\ge\frac3{4(2K+1)}\ge\frac3{4(4\sqrt v+3)}$ and
$\alpha\,\mathrm{TC}(\mu_A)\le\psi(v)I$ with $\psi(v)=\ln\frac{16\sqrt v+12}3\cdot\frac{(1+3v)^{1/2}}{2v}$. Comparing
logarithmic derivatives shows that $\psi$ decreases on $[\frac12,\infty)$, and $\psi(\frac12)<3.25$. Hence
$\alpha\,\mathrm{TC}(\mu_A)\le3.25\,I$ and $V\ln(1/V)\le4\alpha\ln(1/\alpha)\le\sum_j\Inf_j\ln^+(4\rho_j/V)+(1+\ln2+13)I$.
\end{proof}

\begin{lemma}\label{lem:weak-compose}
Let $F=G(f_1,\dots,f_k)$, where $f_1,\dots,f_k$ depend on pairwise disjoint sets of variables, and let $\mu$ be the
product measure with $\mu(y_i=1)=\Pr[f_i=1]$. If $G$ satisfies \eqref{eq:biased} under $\mu$ with $I_\mu(G)$ in place
of $I_\mu(G)-\Var_\mu(G)$ and with constant $C$, and $H^+(f_i)\le D\,I(f_i)$ for every $i$, then
$H^+(F)\le(C+D)\,I(F)$.
\end{lemma}

\begin{proof}
By the composition formula of~\cite[Section~3]{OT13}, the spectral distribution of $F$ is obtained by choosing $A$
with law $P_\mu$ and then, independently for $i\in A$, a nonempty set $T_i$ of variables of $f_i$ with law
$\widehat{f_i}(T_i)^2/\sigma_i^2$, where $\sigma_i^2=\Var(f_i)$. Since $\sum_{T\ne\emptyset}\widehat{f_i}(T)^2=\sigma_i^2$,
this gives
\[
H^+(F)=\sum_{A\ne\emptyset}P_\mu(A)\log\frac{\sigma_A^2}{P_\mu(A)}+\sum_i\frac{\Inf_i^\mu(G)}{\sigma_i^2}H^+(f_i),\qquad
I(F)=\sum_i\frac{\Inf_i^\mu(G)}{\sigma_i^2}I(f_i),
\]
and the claim follows from $\Inf_i^\mu(G)=\frac{\Inf_i^\mu(G)}{\sigma_i^2}\Var(f_i)\le\frac{\Inf_i^\mu(G)}{\sigma_i^2}I(f_i)$.
\end{proof}

\begin{corollary}\label{cor:depth}
Let $f$ be computed by a read-once formula of depth $d$ whose gates are $\mathrm{AND}$, $\mathrm{OR}$ and $\mathrm{XOR}$
gates and threshold gates $T_{t,r}$ of any fan-in, with negations allowed. Then $H(f)\le(38d+\log e)\,I(f)$.
\end{corollary}

\begin{proof}
By \cref{prop:thr-general} and \cref{lem:andor}, every gate satisfies the hypothesis of \cref{lem:weak-compose} under
every product measure with $C=38>26.3/\ln2$; negating an input replaces $q_j$ by $1-q_j$ and changes nothing else. A
literal has $H^+=0$. By induction on the depth, \cref{lem:weak-compose} gives $H^+(f)\le38d\,I(f)$, and
$H(f)=H^+(f)+(1-\Var(f))\log\frac1{1-\Var(f)}\le H^+(f)+\Var(f)\log e\le(38d+\log e)\,I(f)$.
\end{proof}

The inputs of such a formula may also be functions from the classes of \cref{thm:plus-intro}, on disjoint sets of
variables, at the cost of adding their constant. For formulas of bounded depth the constant does not depend on the
fan-ins, while~\cite[Theorem~3]{OT13} gives a constant independent of the depth that grows exponentially with the
fan-ins. The next theorem removes the dependence on $d$.

\subsection*{The strengthened inequality for threshold gates}
The weak form loses only when one input is pivotal with probability close to one. The other inputs are then nearly
deterministic, and the left side of \eqref{eq:biased} can be estimated directly.

\begin{theorem}\label{thm:thr-strong}
Let $1\le t\le r$ and let $\mu$ be a product measure on $\bits^r$ with $\mu(y_j=1)=q_j\in(0,1)$. With
$J=\sum_{|A|\ge2}|A|\,P_\mu(A)$,
\[
\sum_{A\ne\emptyset}P_\mu(A)\ln\frac{\sigma_A^2}{P_\mu(A)}\le263\,J\le526\,\bigl(I_\mu(T_{t,r})-\Var_\mu(T_{t,r})\bigr),
\]
so $T_{t,r}$ satisfies \eqref{eq:biased} with $C=760$ under every product measure.
\end{theorem}

\begin{proof}
Write $N$, $N_{-j}$, $\sigma_j^2$, $\rho_j=\mu(N_{-j}=t-1)$, $I$ and $V$ as in \cref{prop:thr-general}, and put
$J_j=\sigma_j^2\rho_j(1-\rho_j)$. The derivative $d_j=\mathbf 1[N_{-j}=t-1]$ has mean $\rho_j$ and
$P(B\cup\{j\})=\sigma_j^2\widetilde{d_j}(B)^2$ (proof of \cref{lem:thr-rank}), so $P(\{j\})=\sigma_j^2\rho_j^2$,
$\sum_{A\ni j,\,|A|\ge2}P(A)=\sigma_j^2\Var(d_j)=J_j$ and $J=\sum_jJ_j$. Since $J\ge2\sum_{|A|\ge2}P(A)$, we have
$I-V=\sum_{|A|\ge2}(|A|-1)P(A)\ge J/2$, which is the second inequality. The sets of size one contribute
$\sum_j2\sigma_j^2\rho_j^2\ln(1/\rho_j)\le2J$.

If $\rho_j\le\frac9{10}$ for every $j$, then $J\ge I/10$ and \cref{prop:thr-general} gives the bound $263\,J$.
Otherwise let $j_0$ maximise $\rho_j$, so that $\rho_{j_0}=1-\varepsilon$ with $\varepsilon<\frac1{10}$. Negating all
inputs and the output maps $T_{t,r}$ to $T_{r-t+1,r}$ and changes none of the quantities involved, so we may assume
$q:=q_{j_0}\le\frac12$. Let $O=[r]\setminus\{j_0\}$ and $N'=N_{-j_0}$. For $i\in O$ let $\tilde q_i=\min(q_i,1-q_i)$,
let $\bar y_i$ be the more likely value of $y_i$ and $z_i=\mathbf 1[y_i\ne\bar y_i]$; let
$U=\{i\in O:\bar y_i=0\}$, $D=\{i\in O:\bar y_i=1\}$, $Z_U=\sum_{i\in U}z_i$, $Z_D=\sum_{i\in D}z_i$,
$L_U=\sum_{i\in U}\tilde q_i$, $L_D=\sum_{i\in D}\tilde q_i$, $L'=L_U+L_D$, $\Lambda_U=\sum_{i\in U}\sigma_i^2$,
$\Lambda_D=\sum_{i\in D}\sigma_i^2$ and $\Lambda'=\Lambda_U+\Lambda_D$.

\emph{The other inputs are nearly deterministic.} Resampling every coordinate, as in the proof of
\cref{lem:thr-rank} with $\eta=1$, and using $\E[1/(F+1)]\le1/\E F$ as in the proof of \cref{prop:thr-general}, gives
$\mu(N'=k)\le(2/(\pi L'))^{1/2}$ for every $k$, because the number $F$ of free coordinates has mean $2L'$. If $|D|\ne t-1$, then $N'=t-1$ forces some $z_i=1$, so
$\prod_{i\in O}(1-\tilde q_i)\le\varepsilon<\frac1{10}$; as $\ln(1-x)\ge-\frac32x$ for $0\le x\le\frac12$, this gives
$L'>\frac23\ln10>1.53$ and $\rho_{j_0}<0.65$, which is impossible. Hence $|D|=t-1$ and $N'-(t-1)=Z_U-Z_D$. A
configuration with exactly one $z_i=1$ has $N'\ne t-1$, so $\varepsilon\ge L'(1-L')$; since also
$L'\le2/(\pi\rho_{j_0}^2)<0.79$, this forces $L'<0.113$, and then $L'\le1.128\,\varepsilon$ and
$\Lambda'\le4L'\le4.51\,\varepsilon\le0.451$.

\emph{Lower bounds for $J$.} For $i\in U$, $N_{-i}=t-1$ exactly when $y_{j_0}+Z_{U\setminus\{i\}}-Z_D=0$, so
$\rho_i\ge(1-q)(1-L')>0.44$, while $\rho_i\le1-\varepsilon$; by concavity
$\rho_i(1-\rho_i)\ge\min(0.44\cdot0.56,\,\varepsilon(1-\varepsilon))\ge0.9\,\varepsilon$. For $i\in D$, $N_{-i}=t-1$
exactly when $y_{j_0}+Z_U-Z_{D\setminus\{i\}}=1$, so $\rho_i\ge q(1-L')\ge0.887\,q$ and $1-\rho_i\ge(1-q)(1-L')>0.44$,
whence $\rho_i(1-\rho_i)\ge0.49\,q$. Therefore $\sum_{i\in U}J_i\ge0.9\,\varepsilon\Lambda_U$ and
$\sum_{i\in D}J_i\ge0.49\,q\Lambda_D$.

\emph{Sets containing $j_0$.} With $d=d_{j_0}$, the sets $A\ni j_0$ with $|A|\ge2$ contribute
\[
\sigma_{j_0}^2\sum_{B\ne\emptyset}\widetilde d(B)^2\ln\frac{\sigma_B^2}{\widetilde d(B)^2}.
\]
Let $U_k$ be the weight of $d$ at level $k$, $U_d=\Var(d)=\varepsilon(1-\varepsilon)$, and $e_k$ the $k$-th elementary symmetric polynomial of
$(\sigma_i^2)_{i\in O}$, so that $e_k\le\Lambda'^k\le(5.02\,U_d)^k$. Gibbs' inequality on each level against the
weights $\sigma_B^2$, together with $x\ln(a/x)\le a/e$, gives
\[
\sum_{B\ne\emptyset}\widetilde d(B)^2\ln\frac{\sigma_B^2}{\widetilde d(B)^2}\le\sum_kU_k\ln\frac{e_k}{U_k}
\le\sum_k\Bigl(\frac{U_d^k}e+kU_k\ln5.02\Bigr)\le0.41\,U_d+1.62\,I(d),
\]
and $I(d)=\sum_{i\in O}q_i(1-q_i)\,\mu(i\text{ is pivotal for }d)\le\Lambda'/4\le1.26\,U_d$. So these sets contribute
at most $2.5\,\sigma_{j_0}^2U_d=2.5\,J_{j_0}$.

\emph{Sets avoiding $j_0$.} Put $u=\mathbf 1[Z_U-Z_D\ge1]$ and $v=\mathbf 1[Z_D-Z_U\ge1]$. Since
$T_{t,r}=u+y_{j_0}(1-u-v)$, for $\emptyset\ne A\subseteq O$ the coefficient of $\phi_A$ in the $\pm1$ version of
$T_{t,r}$ is $2(1-q)\tilde u(A)-2q\tilde v(A)$, where $\tilde u$ and $\tilde v$ are the coefficients of $u$ and $v$.
Computing in one coordinate of $A$ as in \cref{lem:thr-rank} and using $\E|\phi_i|=\sigma_i$ gives
$|\tilde u(A)|,|\tilde v(A)|\le\sigma_A/2$. Moreover, for $A\subseteq D$ we have
$\tilde u(A)=\E\bigl[\E[u\mid y_D]\,\phi_A\bigr]$ with $0\le\E[u\mid y_D]\le\mu(Z_U\ge1)\le L_U$, so
$|\tilde u(A)|\le L_U\sigma_A$, and likewise $|\tilde v(A)|\le L_D\sigma_A$ for $A\subseteq U$. We use three
elementary facts. For $s>0$ and $a^2,b^2\le s$, the function $\varphi_s(x)=x\ln(s/x)$ satisfies
$\varphi_s((a+b)^2)\le2\varphi_s(a^2)+2\varphi_s(b^2)+2a^2+2b^2$: if $2a^2+2b^2\le s/e$ this follows from the
monotonicity of $\varphi_s$ on $[0,s/e]$ and $(x+y)\ln\frac s{x+y}\le x\ln\frac sx+y\ln\frac sy$, and otherwise
$\varphi_s\le s/e<2a^2+2b^2$. For a family $\mathcal F$ of sets and numbers $0\le x_A\le\lambda\sigma_A^2$ with
$\lambda\le1/e$, the log-sum inequality gives $\sum_{\mathcal F}x_A\ln(\sigma_A^2/x_A)\le\lambda\ln(1/\lambda)
\sum_{\mathcal F}\sigma_A^2$, and $\frac1e\sum_{\mathcal F}\sigma_A^2$ is a bound when only $x_A\le\sigma_A^2$ is
known. Finally, over the sets with $|A|\ge2$ that meet $U$, $\sum\sigma_A^2\le\Lambda_U(e^{\Lambda'}-1)\le
1.57\,\Lambda'\Lambda_U$, and over the sets $A\subseteq D$ with $|A|\ge2$,
$\sum\sigma_A^2\le\frac12\Lambda_D^2e^{\Lambda_D}\le0.79\,\Lambda_D^2$; the same holds with $U$ and $D$ exchanged. We
split the sets $A\subseteq O$ with $|A|\ge2$ according to whether they meet $U$ for the terms with
$a=2(1-q)\tilde u(A)$, and according to whether they meet $D$ for the terms with $b=2q\tilde v(A)$. With $\lambda=1$,
$4L_U^2$, $q^2$ and $4q^2L_D^2$ in the four families, $L_U\le\Lambda_U/2$, $L_D\le\Lambda_D/2$,
$x^2(1+2\ln\frac1x)\le0.91\,x$ for $0<x\le0.23$ and $q(1+2\ln\frac1q)\le1.2$ for $q\le\frac12$, these sets contribute
at most
\[
19.4\,\varepsilon\Lambda_U+2.93\,\varepsilon\Lambda_U+1.7\,q\Lambda_D+0.3\,q\Lambda_D
\le24.9\sum_{i\in U}J_i+4.1\sum_{i\in D}J_i.
\]
Adding the three contributions, the left side is at most $2J+2.5\,J_{j_0}+24.9\,(J-J_{j_0})\le26.9\,J$.
\end{proof}

\begin{corollary}\label{cor:thr-formula}
Let $f$ be computed by a read-once formula whose gates are threshold gates $T_{t,r}$ and $\mathrm{XOR}$ gates of any
fan-in, with negations allowed. Then $H(f)\le760\,I(f)$.
\end{corollary}

\begin{proof}
By \cref{thm:thr-strong} and \cref{lem:andor}, every gate satisfies \eqref{eq:biased} with $C=760$ under every
product measure, and negating an input changes nothing. The composition theorem~\cite[Theorem~1]{OT13}, applied at
each gate from the inputs to the output, gives $H^+(f)\le760\,(I(f)-\Var(f))$, and
$H(f)\le H^+(f)+\Var(f)\log e\le760\,I(f)$.
\end{proof}

The inputs of such a formula may again be functions from the classes of \cref{thm:plus-intro}, and gates of bounded
arity may be added, with the constants of \cref{cor:compose-intro}. For depth $d\le19$, \cref{cor:depth} gives the
smaller constant $38d+\log e$, and the numerical ratio in \cref{sec:numerics} is below $2+4\log e$ for every threshold
gate that was computed.


\begin{proof}[Proof of \cref{cor:compose-intro}]
By~\cite[Theorem~2]{OT13}, every gate of arity at most $a$ satisfies the biased form of \eqref{eq:feiplus} with a
constant $C_a=2^{O(a)}$, for every product measure on its inputs, and by \cref{lem:andor} so do $\mathrm{AND}$,
$\mathrm{OR}$ and $\mathrm{XOR}$ gates of any fan-in, with constant $7.78$; by \cref{thm:threshold} so does a threshold gate whose inputs have equal biases, with constant
$66.4$, and by \cref{thm:thr-strong} so does every threshold gate, with constant $760$. By \cref{thm:plus-intro},
each $f_l$ satisfies
\eqref{eq:feiplus} with the constant of its class, which is at least $50$, and so does every gate of the last kind,
because its inputs are independent balanced bits and the product measure that~\cite[Theorem~1]{OT13} puts on them is
the uniform one. The composition theorem~\cite[Theorem~1]{OT13}, applied at each gate of $F$ from the inputs to the
output, shows that $f$ satisfies \eqref{eq:feiplus} with constant $C$, and \cite[Proposition~2.1]{OT13} gives the
bound on $H(f)$.
\end{proof}

\begin{remark}
Only the uniform form of \eqref{eq:feiplus} is needed for the inputs $f_l$, because they read the uniformly
distributed variables directly, and for gates with balanced inputs. Using read-$k$ DNF formulas or regular
functions as gates with biased inputs would require the biased form of \eqref{eq:feiplus} for these classes, which
the proofs above do not give, except for $\mathrm{AND}$ and $\mathrm{OR}$ (\cref{lem:andor}) and for threshold gates
(\cref{thm:threshold,thm:thr-strong}).
\end{remark}

\section{Numerical examples}\label{sec:numerics}

All quantities below were computed exactly from full truth tables by the Walsh--Hadamard transform. For read-$k$
DNF formulas the table lists $H/I$ and the right side of \eqref{eq:gibbs}, evaluated exactly with the reference
measure of \cref{thm:readk-main} and divided by $I$; this is the quantity that the proof of \cref{thm:readk-main}
bounds from above.

\begin{center}
\begin{tabular}{lcccc}
\toprule
formula & $n$ & $k$ & $H/I$ & Gibbs bound $/I$ \\
\midrule
tribes, $3$ terms of width $5$ & 15 & 1 & 3.632 & 3.632\\
stars of a random $3$-regular graph on $12$ vertices & 18 & 2 & 4.128 & 4.366\\
stars of a random $3$-regular graph on $16$ vertices & 24 & 2 & 4.487 & 4.724\\
stars of a random $4$-regular graph on $12$ vertices & 24 & 2 & 3.875 & 4.004\\
random monotone read-$3$ formulas, larger of $2$ & 20 & 3 & 3.942 & 4.145\\
random read-$k$ formulas with negations, $192$ samples & $\le16$ & $\le5$ & 3.781 & 4.524\\
\bottomrule
\end{tabular}
\end{center}

\noindent Here the star function of a graph $G$ has one variable per edge and is true when some vertex has all its
incident edges present; it is computed by a read-$2$ DNF formula with one term per vertex. The last row lists the
largest values over the $192$ formulas; the two maxima need not be attained by the same formula. On all these
examples the Gibbs bound exceeds the entropy by at most $1.32\,I$, and the largest ratio is $H/I=4.49$, for the star
function of a $3$-regular graph. For the reference measure of \cref{thm:readk-log} with $\beta=3$, the upper bound on
the right side of \eqref{eq:gibbs} given by the runs constructed in its proof, computed from the exact spectral
distribution of $28$ random read-$k$ formulas with $k\le10$ and $12$ or $14$ variables, lies between $4.5\,I$ and
$8.4\,I$ and grows slowly with $k$, while $H/I\le3.91$ on the same formulas; the probabilities of these runs sum to
at most $0.12$ over all sets. For regular functions, $590$ random linear threshold functions and $440$ random
regular functions obtained as up-sets of shifted families, all with at most $13$ variables, have $H/I\le3.94$.
Intersections, unions, parities and majorities of two to four linear threshold functions on at most $12$
variables, with independent random weights or with weights in opposite orders, have $H/I\le3.86$; for these the
quantity $\E_S\sum_{i\in S}\min_l\ln^+(r_l(i)/|S|)$ in the proof of \cref{thm:multi} is at most $0.24\,I(f)$, and
the expectation of $F_l$ under each piece $\nu_{i,l}$ is at most $0.74$, against the bounds $4\ln m+4.93$ and $L_1$.
For $75$ read-$k$ unions of threshold functions with $k\le4$, up to $12$ threshold functions and at most $14$
variables, $36$ of them built from rarely true threshold functions on $5$ to $9$ variables, $H/I\le4.0$, the ratio of
the two sides of \cref{lem:upath} is at most $0.96$, the expectation of $F^j_{l_s}$ under each piece is at most $0.73$,
the quantity bounded by $K$ in the proof of \cref{thm:union} is at most $0.30$, and the upper bound on $H(f)$ given by
the choice of $L_A$ in that proof, computed from the exact spectral distribution, lies between $H(f)$ and $37\,I(f)$.
For the reference measure of \cref{thm:union-log} with $a=1.44$ and $\beta=10$, the bound given by the runs
constructed in its proof, on $13$ random read-$k$ unions of threshold functions with $k\le5$ and $10$ or $12$
variables, lies between $4.9\,I$ and $7.9\,I$, and the run probabilities sum to at most $0.17$ over all sets.
For the reference measure of \cref{thm:union-multi}, on $22$ random read-$k$ unions of intersections of two or three
threshold functions with $k\le3$ and $10$ or $12$ variables, the bound lies between $4.0\,I$ and $7.2\,I$, the run
probabilities sum to at most $0.16$, and the per-set inequality of the proof holds on every set with a margin of at
least $13$ bits; on $15$ read-$k$ unions of random functions of $w\le5$ variables each ($k\le3$, $10$ or $12$
variables), viewed as functions of dictators, the bound lies between $4.0\,I$ and $7.6\,I$, the run probabilities sum
to at most $0.21$, and the margin is at least $14$ bits.

For the strengthened inequality, the largest values of $H^+(f)/(I(f)-\Var(f))$ observed were $6.35$ on random
read-$k$ DNF formulas with $k\le6$ and at most $14$ variables, $6.11$ on $320$ random regular and linear threshold
functions with at most $12$ variables, $5.68$ on intersections, unions, parities and majorities of two to four linear
threshold functions on the same variables, and $6.38$ on conjunctions, disjunctions, parities and majorities of such
functions on disjoint sets of variables. In \cref{lem:mass} the largest observed value of
$W(f)\log(1/W(f))/J(f)$, over all functions of at most $4$ variables, random functions of up to $10$ variables and
perturbations of dictators, was $1.59$.

For threshold gates with inputs of equal bias every quantity in \cref{thm:threshold} was computed with $50$-digit
arithmetic from the formula $W_k=\binom rk(q(1-q))^k\bigl(\E\,\Delta^kG(c)\bigr)^2$, where $G(c)=1-2\cdot\mathbf 1[c\ge t]$,
$\Delta$ is the forward difference and $c$ has the law $\Bin(r-k,q)$: for all $r\le12$, all $1\le t\le r$ and
$q\in\{10^{-e},1-10^{-e}:e=1,2,4,8,16,30\}\cup\{0.1,0.3,0.5,0.7,0.9\}$, and for $150$ random triples with
$13\le r\le40$. The left side of \cref{thm:threshold} divided by $I$ is at most $2.94$ (the bound is $11.49$); in
\cref{lem:thr-rank} the ratio is at most $1.14$ (bound $4.85$); in \cref{lem:thr-credit} the quantity
$(V\ln(1/V)-I\ln^+(4\rho/V))/I$ is at most $0.68$ (bound $4.64$) and $\alpha\,\mathrm{TC}(\mu_A)/I$ is at most $0.25$
(bound $2/e$). The ratio of the two sides of \eqref{eq:biased} is at most $7.7708$, the value $2+4\log e$ of
\cref{lem:andor}, attained by $\mathrm{OR}_2$ as $q\to0$. For linear threshold functions with at most $13$ variables
under $3{,}000$ random product measures with biases down to $10^{-12}$, $\alpha\,\mathrm{TC}(\mu_A)\le0.24\,I_\mu$ as
well, while for the disjunctions of \cref{rem:threshold} the ratio grows from $0.23$ to $0.39$ as $s$ goes from $3$
to $7$. For $3{,}000$ further linear threshold functions with at most $12$ variables, with weights in $\{1,2,3\}$,
powers of two up to $2^6$, or exponentially distributed, $\alpha\,\mathrm{TC}(\mu_A)\le0.31\,I_\mu$; for these the
bound $\mathrm{TC}(\mu_A)\le\ln(\chi^*/\alpha)$ given by the exponential tilt, where $\chi^*$ is the Chernoff bound for
$\alpha$, is far from sufficient: $\alpha\ln(\chi^*/\alpha)/I_\mu$ reaches $1.9\cdot10^3$ on an event of probability
about $10^{-16}$. Over $1{,}500$ random linear threshold functions with at most $9$ variables, and $60$ searches that
maximise the ratio of the two sides of \eqref{eq:biased} over the weights, the threshold and the biases of a linear
threshold function with at most $8$ variables, that ratio never exceeded $2+4\log e$; nor did it for one heavy input
and up to $32$ light inputs of a common bias, computed exactly by level over a grid of biases and thresholds, where the
largest value is always attained at the conjunction or disjunction corner. For $1{,}200$ threshold functions with at most $8$ variables and unequal biases down to $10^{-14}$, the
per-level Gibbs bound in the proof of \cref{prop:thr-general}, computed exactly, lies between the left side and
$2.1\,I_\mu$ (the bound is $26.3\,I_\mu$), the rank terms are at most $0.79\,I_\mu$ and $\alpha\,\mathrm{TC}(\mu_A)$
is at most $0.17\,I_\mu$; and for $1{,}500$ threshold functions with at most $9$ variables the ratio of the two sides
of \eqref{eq:biased} is again at most $2+4\log e$. For \cref{thm:thr-strong}, $1{,}600$ threshold functions with at
most $10$ variables were computed in $50$-digit arithmetic, half of them built so that one input is pivotal with
probability above $\frac9{10}$. In that case $L'\le1.05\,\varepsilon$ (the proof uses $1.128\,\varepsilon$),
$\Lambda'\le4.07\,\varepsilon$ (it uses $4.51\,\varepsilon$), the sets containing $j_0$ contribute at most
$1.57\,J_{j_0}$ (it uses $2.5\,J_{j_0}$), the sets avoiding $j_0$ at most $6\%$ of the bound used for them, and the
left side is at most $2.7\,J$ (it uses $26.9\,J$); in the other case the left side is at most $2.1\,I_\mu$.

\section{Concluding remarks}

The proofs for read-once formulas~\cite{OT13}, for regular functions (\cref{thm:regular,thm:multi}) and for
read-$k$ DNF formulas (\cref{thm:readk}) use unrelated mechanisms: composition, a rank order with the
Littlewood--Offord inequality, and activation of witnessing terms. \Cref{thm:multi} shows that the rank mechanism
tolerates several unrelated orders at a logarithmic cost, \cref{thm:union} that the last two mechanisms combine
when the regular functions overlap boundedly, and \cref{thm:union-multi} that the three combine. The rank order suits threshold-like functions and has no analogue
for read-once formulas such as recursive majority, whose automorphism group is transitive on the variables. For the
majority function each variable lies in exponentially many minterms, and the contributions of the corresponding
subcubes to the derivative superpose coherently, so \cref{lem:witness} gives nothing.

For read-$k$ DNF formulas the dependence on $k$ in \cref{thm:readk} is logarithmic, as for the min-entropy
bound of~\cite{ACKSW21}. A bound independent of $k$ would be the conjecture for all DNF formulas, since a formula
with $m$ terms is read-$m$; on the examples we computed, including sunflower formulas, vertex stars of complete $2$-
and $3$-uniform hypergraphs, graph formulas $\bigvee_{uv\in E}(x_u\wedge x_v)$ of random $k$-regular graphs with
$k\le10$ and random read-$k$ formulas, $H/I$ stays below $4.5$ and does not grow with $k$. For unions of regular
functions (\cref{thm:union-log}) the logarithm enters in the same two ways: through the choice among at most $k$
overlapping functions, and through the Cauchy--Schwarz step, whose factor $k$ costs only $\log k$ because the rank
tails and the level tails of the pieces decay exponentially. A second
question is whether the biased form \eqref{eq:biased} holds for read-$k$ DNF formulas and for regular functions
under all product measures; together with~\cite{OT13} this would allow such functions as gates of read-once
formulas, not only as inputs, and for threshold functions it would give the conjecture for all read-once threshold
formulas. Every term on the left of \eqref{eq:biased} is nonnegative: computing $\E[g\phi_A]$ one coordinate at a time
gives $P_\mu(A)=\sigma_A^2\,4^{-|A|}\bigl(\E\,\partial_Ag\bigr)^2$, where
$\partial_Ag(x)=\sum_{B\subseteq A}(-1)^{|A\setminus B|}g(x^{A\leftarrow1_B})$ is the mixed difference in the
coordinates of $A$ and the expectation is over the other coordinates, so the left side of \eqref{eq:biased} equals
$\sum_{A\ne\emptyset}P_\mu(A)\log\bigl(4^{|A|}/(\E\,\partial_Ag)^2\bigr)$ and involves the biases of the coordinates in
$A$ only through the weight $\sigma_A^2$. The numerical evidence points to the constant of \cref{lem:andor}. For $60$ threshold functions with at
most $11$ variables and $40$ read-$k$ DNF formulas with $k\le3$ and at most $12$ variables, with the product measure
chosen by local search, the largest ratio of the two sides of \eqref{eq:biased}, recomputed with $50$-digit
arithmetic, is below $2+4\log e$ and approaches it only at extreme biases, where the function behaves like
$\mathrm{AND}_2$. The same holds, in exact arithmetic, for every symmetric function of at most $8$ variables and
every symmetric threshold function of at most $30$ variables under product measures with equal biases, at
fifteen biases between $10^{-8}$ and $1-10^{-8}$. For symmetric threshold functions the ratio is
bounded by \cref{thm:threshold} under equal biases and by \cref{thm:thr-strong} in general, with a larger constant. Its proof isolates what a general proof needs: after the rank step, the only
obstacle is the normalisation $V\ln(1/V)$, which the chain rule splits into the marginal divergences of the
conditional law $\mu_A$, all of them repaid by the factors $\sigma_A^2$, and the total correlation of $\mu_A$. For
threshold functions under all product measures the total correlation is controlled by the tilt, which gives
\cref{prop:thr-general}; when one input is pivotal with probability close to one, the other inputs are nearly
deterministic and \cref{thm:thr-strong} follows by direct estimates. For general linear threshold functions the numerical bound $\alpha\,\mathrm{TC}(\mu_A)\le0.31\,I_\mu$ of
\cref{sec:numerics} is what a proof would have to establish, and the exponential tilt that gives it for symmetric
thresholds does not: with unequal weights the Chernoff bound can exceed a rare event's probability by a factor that
is not controlled by its influence. The inequality \eqref{eq:biased} itself held with the constant $2+4\log e$ of
\cref{lem:andor} for every linear threshold function that was computed, and a proof would extend
\cref{cor:thr-formula} to read-once formulas over linear threshold gates of any fan-in.

A route that avoids the total correlation is to divide each set among its elements. For every Boolean function and
every product measure,
\[
\sum_{|A|\ge2}P(A)\ln\frac{\sigma_A^2}{P(A)}=\sum_j\sigma_j^2\,\Psi(d_j),\qquad
\Psi(d)=\sum_{B\ne\emptyset}\frac{\widetilde d(B)^2}{|B|+1}\ln\frac{\sigma_B^2}{\widetilde d(B)^2},
\]
where $d_j$ is the derivative in the $j$-th coordinate. For a function whose derivatives do not change sign, such as
a linear threshold function, the sets of size one contribute at most $2J$, so a bound $\Psi(d_j)\le C\Var(d_j)$ for
every $j$ gives \eqref{eq:biased} with constant $2(2+C)/\ln2$. For a threshold gate $d_j$ is the indicator of $\{N_{-j}=t-1\}$, whose
coefficient on a set $B$ of size $b$ is $\prod_{i\in B}\sqrt{q_i(1-q_i)}$ times the $b$-th backward difference, at
$t-1$, of the law of the number of ones outside $B$. For the conjunction of $m$ inputs of bias $q$, the slice
$N=m$, the ratio $\Psi/\Var$ tends to $\ln4+2q\ln(1/q)/(1-q)$ as $m(1-q)\to\infty$, and every slice that was
computed, with up to $256$ inputs of equal bias or up to $8$ inputs of unequal biases, satisfies
$\Psi<(2+\ln4)\Var$. For the conjunction this holds under every product measure: there
$\widetilde d(B)^2=\rho\prod_{i\in B}(1-q_i)\prod_{i\notin B}q_i$, and $q\ln(1/q)\le1-q$ bounds the part of $\Psi/\Var$
coming from the inputs outside $B$ by $2$, so $2+\ln4$ is the supremum over conjunctions. If $\Psi<(2+\ln4)\Var$ held
for all slices, \eqref{eq:biased} would follow for threshold gates with constant $2(4+\ln4)/\ln2<15.6$. When $\rho_j\ge\frac12$, the anti-concentration bound used in the proof of
\cref{thm:thr-strong} forces the other inputs to have $\sigma^2$-mass at most $930\Var(d_j)$, and Gibbs' inequality
on each level then gives $\Psi(d_j)\le7.04\Var(d_j)$; the open case is a slice far in the tail of the law of $N_{-j}$. For a linear threshold function the bound fails
coordinate by coordinate: the derivative of a heavy input can be a threshold function of the light inputs, and then
$\Psi/\Var$ grows like the logarithm of the number of light inputs, although the whole left side stays small.
Other divisions of the sets among their elements do better: for $\mathbf 1[m\,y_0+y_1+\cdots+y_m\ge m/2]$ with inputs
of bias $\frac12$ the equal division gives the heavy input the ratio $4.12$ at $m=512$, while giving each set to its
lightest elements keeps every coordinate at $1.88$, and the best division, found by linear programming, keeps every
coordinate below $1.84$ for every computed linear threshold function with at most $7$ inputs. What is missing for
linear threshold functions is an explicit division with this property.

\section*{Funding}
The research presented in this paper was supported by the European Research Council (ERC) under the European
Union's Horizon Europe research and innovation programme (grant agreement No.~101041711), by the Simons Foundation
as part of the Collaboration on the Mathematical and Scientific Foundations of Deep Learning, by Heights Labs, by
Convex Nexus Capital, by the Israel Science Foundation (grant number 2258/19), and by the Israel Science Foundation
(ISF Grant 4101/25). Computational resources were provided in part by Google Cloud through its education credits
program (credit 530550103).

\section*{Declaration of competing interest}
The author holds equity in Heights Labs and in Convex Nexus Capital and has a management role at Convex Nexus
Capital. These entities are included in the funding declaration above.

\section*{Declaration of generative AI and AI-assisted technologies}
The theorems, proofs and numerical computations in this paper were produced by the AI system Claude (Anthropic),
working within the author's math-brain project. The case of linear threshold functions in \cref{thm:regular} was
produced earlier in the same project by Claude agents, and the case of Boolean functions of several threshold
functions in \cref{thm:multi} was also obtained there by another group of Claude agents, with a different argument
(interleaving the orders) and a smaller constant for large $m$. The author takes full responsibility for the content
of the paper.

\begin{thebibliography}{99}

\bibitem{ACKSW21} S.~Arunachalam, S.~Chakraborty, M.~Kouck\'y, N.~Saurabh and R.~de~Wolf, Improved bounds on
Fourier entropy and min-entropy, \emph{37th International Symposium on Theoretical Aspects of Computer Science
(STACS 2020)}, LIPIcs 154, 45:1--45:19; \emph{ACM Trans. Comput. Theory} 13 (2021).

\bibitem{Ber67} A.~J.~Bernstein, Maximally connected arrays on the $n$-cube, \emph{SIAM J. Appl. Math.} 15 (1967),
1485--1489.

\bibitem{CKKLS18} S.~Chakraborty, S.~Karmalkar, S.~Kundu, S.~V.~Lokam and N.~Saurabh, Fourier entropy-influence
conjecture for random linear threshold functions, \emph{LATIN 2018}, Lecture Notes in Comput. Sci. 10807,
Springer, 2018, 275--289.

\bibitem{CKLS16} S.~Chakraborty, R.~Kulkarni, S.~V.~Lokam and N.~Saurabh, Upper bounds on Fourier entropy,
\emph{Theoret. Comput. Sci.} 654 (2016), 92--112.

\bibitem{Dar64} J.~N.~Darroch, On the distribution of the number of successes in independent trials, \emph{Ann. Math.
Statist.} 35 (1964), 1317--1321.

\bibitem{FK96} E.~Friedgut and G.~Kalai, Every monotone graph property has a sharp threshold, \emph{Proc. Amer.
Math. Soc.} 124 (1996), 2993--3002.

\bibitem{GMP26} M.~J.~Gonz\'alez, P.~MacManus and M.~C.~Pereyra, Further evidence towards the Fourier
entropy-influence conjecture, arXiv:2606.00246 (2026).

\bibitem{Han25} X.~Han, A new bound for the Fourier-entropy-influence conjecture, \emph{Combinatorica} 45 (2025).

\bibitem{Har64} L.~H.~Harper, Optimal assignments of numbers to vertices, \emph{J. Soc. Indust. Appl. Math.} 12
(1964), 131--135.

\bibitem{Har76} S.~Hart, A note on the edges of the $n$-cube, \emph{Discrete Math.} 14 (1976), 157--163.

\bibitem{Hoe56} W.~Hoeffding, On the distribution of the number of successes in independent trials, \emph{Ann. Math.
Statist.} 27 (1956), 713--721.

\bibitem{Hod17} R.~Hod, Improved lower bound for the Fourier entropy/influence conjecture via lexicographic
functions, arXiv:1711.00762 (2017).

\bibitem{KB80} R.~Kaas and J.~M.~Buhrman, Mean, median and mode in binomial distributions, \emph{Statist. Neerlandica}
34 (1980), 13--18.

\bibitem{KKLMS20} E.~Kelman, G.~Kindler, N.~Lifshitz, D.~Minzer and M.~Safra, Towards a proof of the
Fourier--entropy conjecture?, \emph{Geom. Funct. Anal.} 30 (2020), 1097--1138.

\bibitem{KLW10} A.~Klivans, H.~Lee and A.~Wan, Mansour's conjecture is true for random DNF formulas,
\emph{Proceedings of the 23rd Conference on Learning Theory (COLT 2010)}, 368--380.

\bibitem{KMS12} N.~Keller, E.~Mossel and T.~Schlank, A note on the entropy/influence conjecture, \emph{Discrete
Math.} 312 (2012), 3364--3372.

\bibitem{Man95} Y.~Mansour, An $O(n^{\log\log n})$ learning algorithm for DNF under the uniform distribution,
\emph{J. Comput. System Sci.} 50 (1995), 543--550.

\bibitem{MS77} F.~J.~MacWilliams and N.~J.~A.~Sloane, \emph{The Theory of Error-Correcting Codes}, North-Holland,
Amsterdam, 1977.

\bibitem{Mur71} S.~Muroga, \emph{Threshold Logic and its Applications}, Wiley-Interscience, New York, 1971.

\bibitem{Num26} Numaro, Certificate archive: improved lower bound $C_{71}>6.521845710923046575$ for the Fourier
entropy-influence constant ($n=18$ explicit construction), Zenodo, 2026, doi:10.5281/zenodo.21497769.

\bibitem{OT13} R.~O'Donnell and L.-Y.~Tan, A composition theorem for the Fourier entropy-influence conjecture,
\emph{ICALP 2013}, Lecture Notes in Comput. Sci. 7965, Springer, 2013, 780--791.

\bibitem{OWZ11} R.~O'Donnell, J.~Wright and Y.~Zhou, The Fourier entropy-influence conjecture for certain classes
of Boolean functions, \emph{ICALP 2011}, Lecture Notes in Comput. Sci. 6755, Springer, 2011, 330--341.

\bibitem{Pot26} V.~N.~Potapov, On the Fourier entropy-influence conjecture for Boolean plateaued functions,
arXiv:2609.21563, 2026.

\bibitem{Sha18} G.~Shalev, On the Fourier entropy influence conjecture for extremal classes, arXiv:1806.03646
(2018).

\bibitem{Spe28} E.~Sperner, Ein Satz \"uber Untermengen einer endlichen Menge, \emph{Math. Z.} 27 (1928),
544--548.

\bibitem{WWW14} A.~Wan, J.~Wright and C.~Wu, Decision trees, protocols and the entropy-influence conjecture,
\emph{Proceedings of the 5th Conference on Innovations in Theoretical Computer Science (ITCS 2014)}, 67--80.

\end{thebibliography}

\end{document}
